\begin{filecontents*}[overwrite]{arc_references.bib}
@article{memgpt,
  author = {Packer, Charles and Wooders, Sarah and Lin, Kevin and Fang, Vivian and Patil, Shishir G. and Stoica, Ion and Gonzalez, Joseph E.},
  title = {{MemGPT}: Towards {LLM}s as Operating Systems},
  journal = {arXiv preprint arXiv:2310.08560},
  year = {2023},
  url = {https://arxiv.org/abs/2310.08560}
}
@inproceedings{hiagent,
  author = {Hu, Mengkang and Chen, Tianxing and Chen, Qiguang and Mu, Yao and Shao, Wenqi and Luo, Ping},
  title = {{H}i{A}gent: Hierarchical Working Memory Management for Solving Long-Horizon Agent Tasks with Large Language Model},
  booktitle = {Proceedings of the 63rd Annual Meeting of the Association for Computational Linguistics (Volume 1: Long Papers)},
  pages = {32779--32798},
  year = {2025},
  doi = {10.18653/v1/2025.acl-long.1575},
  url = {https://aclanthology.org/2025.acl-long.1575/}
}
@inproceedings{readagent,
  author = {Lee, Kuang-Huei and Chen, Xinyun and Furuta, Hiroki and Canny, John and Fischer, Ian},
  title = {A Human-Inspired Reading Agent with Gist Memory of Very Long Contexts},
  booktitle = {Proceedings of the 41st International Conference on Machine Learning},
  series = {Proceedings of Machine Learning Research},
  volume = {235},
  pages = {26396--26415},
  year = {2024},
  url = {https://proceedings.mlr.press/v235/lee24c.html}
}
@inproceedings{amem,
  author = {Xu, Wujiang and Liang, Zujie and Mei, Kai and Gao, Hang and Tan, Juntao and Zhang, Yongfeng},
  title = {{A-Mem}: Agentic Memory for {LLM} Agents},
  booktitle = {Advances in Neural Information Processing Systems},
  volume = {38},
  pages = {17577--17604},
  year = {2025},
  doi = {10.52202/085713-0593},
  url = {https://proceedings.neurips.cc/paper_files/paper/2025/hash/19909c36f51abc4856b4560aff3d36d6-Abstract-Conference.html}
}
@inproceedings{acon,
  author = {Kang, Minki and Chen, Wei-Ning and Han, Dongge and Inan, Huseyin and Wutschitz, Lukas and Chen, Yanzhi and Sim, Robert and Rajmohan, Saravan},
  title = {{ACON}: Optimizing Context Compression for Long-horizon {LLM} Agents},
  booktitle = {ICML 2026},
  month = jun,
  year = {2026},
  url = {https://www.microsoft.com/en-us/research/publication/acon-optimizing-context-compression-for-long-horizon-llm-agents/}
}
@inproceedings{agentfold,
  author = {Ye, Rui and Zhang, Zhongwang and Li, Kuan and Yin, Huifeng and Tao, Zhengwei and Zhao, Yida and Su, Liangcai and Zhang, Liwen and Qiao, Zile and Wang, Xinyu and Xie, Pengjun and Huang, Fei and Zhou, Jingren and Chen, Siheng and Jiang, Yong},
  title = {{AgentFold}: Long-Horizon Web Agents with Proactive Context Folding},
  booktitle = {International Conference on Learning Representations},
  pages = {154754--154776},
  year = {2026},
  url = {https://proceedings.iclr.cc/paper_files/paper/2026/hash/fa948624dfde013671e72c1a7ca4aebc-Abstract-Conference.html}
}
@inproceedings{mem1,
  author = {Zhou, Zijian and Qu, Ao and Wu, Zhaoxuan and Kim, Sunghwan and Prakash, Alok and Rus, Daniela and Low, Bryan Kian Hsiang and Liang, Paul},
  title = {{MEM1}: Learning to Synergize Memory and Reasoning for Efficient Long-Horizon Agents},
  booktitle = {International Conference on Learning Representations},
  pages = {58413--58438},
  year = {2026},
  url = {https://proceedings.iclr.cc/paper_files/paper/2026/hash/5fc8b3bdfbb9167b5144df5d3fae4616-Abstract-Conference.html}
}
@misc{contextweaver,
  author = {Wu, Yating and Zhang, Yuhao and Ghosh, Sayan and Basu, Sourya and Deoras, Anoop and Huan, Jun and Gupta, Gaurav},
  title = {{ContextWeaver}: Selective and Dependency-Structured Memory Construction for {LLM} Agents},
  year = {2026},
  eprint = {2604.23069},
  archivePrefix = {arXiv},
  primaryClass = {cs.CL},
  url = {https://arxiv.org/abs/2604.23069}
}
@inproceedings{react,
  author = {Yao, Shunyu and Zhao, Jeffrey and Yu, Dian and Du, Nan and Shafran, Izhak and Narasimhan, Karthik R. and Cao, Yuan},
  title = {{ReAct}: Synergizing Reasoning and Acting in Language Models},
  booktitle = {The Eleventh International Conference on Learning Representations},
  year = {2023},
  url = {https://openreview.net/forum?id=WE_vluYUL-X}
}
@inproceedings{agentspec,
  author = {Wang, Haoyu and Poskitt, Christopher M. and Sun, Jun},
  title = {{AgentSpec}: Customizable Runtime Enforcement for Safe and Reliable {LLM} Agents},
  booktitle = {2026 IEEE/ACM 48th International Conference on Software Engineering (ICSE)},
  year = {2026},
  url = {https://conf.researchr.org/details/icse-2026/icse-2026-research-track/29/AgentSpec-Customizable-Runtime-Enforcement-for-Safe-and-Reliable-LLM-Agents}
}
@misc{forge,
  author = {Palumbo, Nils and Choudhary, Sarthak and Choi, Jihye and Amir, Guy and Chalasani, Prasad and Jha, Somesh},
  title = {Formal Policy Enforcement for Real-World Agentic Systems},
  year = {2026},
  eprint = {2602.16708},
  archivePrefix = {arXiv},
  primaryClass = {cs.CR},
  url = {https://arxiv.org/abs/2602.16708}
}
@inproceedings{guardagent,
  author = {Xiang, Zhen and Zheng, Linzhi and Li, Yanjie and Hong, Junyuan and Li, Qinbin and Xie, Han and Zhang, Jiawei and Xiong, Zidi and Xie, Chulin and Yang, Carl and Song, Dawn and Li, Bo},
  title = {{GuardAgent}: Safeguard {LLM} Agents via Knowledge-Enabled Reasoning},
  booktitle = {Proceedings of the 42nd International Conference on Machine Learning},
  series = {Proceedings of Machine Learning Research},
  volume = {267},
  pages = {68316--68342},
  year = {2025},
  url = {https://proceedings.mlr.press/v267/xiang25a.html}
}
@misc{groundedcontinuation,
  author = {He, Qisong and Dong, Yi and Huang, Xiaowei},
  title = {Grounded Continuation: A Linear-Time Runtime Verifier for {LLM} Conversations},
  year = {2026},
  eprint = {2605.14175},
  archivePrefix = {arXiv},
  primaryClass = {cs.AI},
  url = {https://arxiv.org/abs/2605.14175}
}
@article{sagallm,
  author = {Chang, Edward Y. and Geng, Longling},
  title = {{SagaLLM}: Context Management, Validation, and Transaction Guarantees for Multi-Agent {LLM} Planning},
  journal = {Proceedings of the VLDB Endowment},
  volume = {18},
  number = {12},
  pages = {4874--4886},
  month = aug,
  year = {2025},
  doi = {10.14778/3750601.3750611},
  url = {https://www.vldb.org/pvldb/vol18/p4874-chang.pdf}
}
@misc{atomix,
  author = {Mohammadi, Bardia and Potamitis, Nearchos and Klein, Lars and Arora, Akhil and Bindschaedler, Laurent},
  title = {{Atomix}: Timely, Transactional Tool Use for Reliable Agentic Workflows},
  month = may,
  year = {2026},
  eprint = {2602.14849},
  archivePrefix = {arXiv},
  primaryClass = {cs.LG},
  url = {https://arxiv.org/abs/2602.14849}
}
@misc{sbus,
  author = {Khan, Sajjad},
  title = {{S-Bus}: Automatic Read-Set Reconstruction for Multi-Agent {LLM} State Coordination},
  month = may,
  year = {2026},
  eprint = {2605.17076},
  archivePrefix = {arXiv},
  primaryClass = {cs.LG},
  url = {https://arxiv.org/abs/2605.17076}
}
@misc{committimeauth,
  author = {Santos-Grueiro, Igor},
  title = {Temporary Authority, Permanent Effects: Commit-Time Authorization for {LLM} Agents},
  month = jul,
  year = {2026},
  eprint = {2607.10487},
  archivePrefix = {arXiv},
  primaryClass = {cs.CR},
  url = {https://arxiv.org/abs/2607.10487}
}
@article{meyer1992contract,
  author = {Meyer, Bertrand},
  title = {Applying {``Design by Contract''}},
  journal = {Computer},
  volume = {25},
  number = {10},
  pages = {40--51},
  month = oct,
  year = {1992},
  doi = {10.1109/2.161279},
  url = {https://doi.org/10.1109/2.161279}
}
@article{schneider2000enforcement,
  author = {Schneider, Fred B.},
  title = {Enforceable Security Policies},
  journal = {ACM Transactions on Information and System Security},
  volume = {3},
  number = {1},
  pages = {30--50},
  month = feb,
  year = {2000},
  doi = {10.1145/353323.353382},
  url = {https://doi.org/10.1145/353323.353382}
}
@inproceedings{havelund2002monitors,
  author = {Havelund, Klaus and Ro{\c{s}}u, Grigore},
  title = {Synthesizing Monitors for Safety Properties},
  booktitle = {Tools and Algorithms for the Construction and Analysis of Systems},
  series = {Lecture Notes in Computer Science},
  volume = {2280},
  pages = {342--356},
  year = {2002},
  doi = {10.1007/3-540-46002-0_24},
  url = {https://doi.org/10.1007/3-540-46002-0_24}
}
@article{kung1981occ,
  author = {Kung, H. T. and Robinson, John T.},
  title = {On Optimistic Methods for Concurrency Control},
  journal = {ACM Transactions on Database Systems},
  volume = {6},
  number = {2},
  pages = {213--226},
  month = jun,
  year = {1981},
  doi = {10.1145/319566.319567},
  url = {https://doi.org/10.1145/319566.319567}
}
@article{bernstein1983mvcc,
  author = {Bernstein, Philip A. and Goodman, Nathan},
  title = {Multiversion Concurrency Control---Theory and Algorithms},
  journal = {ACM Transactions on Database Systems},
  volume = {8},
  number = {4},
  pages = {465--483},
  month = dec,
  year = {1983},
  doi = {10.1145/319996.319998},
  url = {https://doi.org/10.1145/319996.319998}
}
@inproceedings{buneman2001provenance,
  author = {Buneman, Peter and Khanna, Sanjeev and Tan, Wang-Chiew},
  title = {Why and Where: A Characterization of Data Provenance},
  booktitle = {Database Theory---ICDT 2001},
  series = {Lecture Notes in Computer Science},
  volume = {1973},
  pages = {316--330},
  year = {2001},
  doi = {10.1007/3-540-44503-X_20},
  url = {https://doi.org/10.1007/3-540-44503-X_20}
}
@inproceedings{taubench,
  author = {Shunyu Yao and Noah Shinn and Pedram Razavi and Karthik Narasimhan},
  title = {{$\tau$}-bench: A Benchmark for Tool-Agent-User Interaction in Real-World Domains},
  booktitle = {International Conference on Learning Representations},
  year = {2025},
  url = {https://proceedings.iclr.cc/paper_files/paper/2025/hash/1b126cc38b8638e07bef37e7b2bb72bf-Abstract-Conference.html}
}
@misc{deepswebenchmark,
  author = {Wenqi Huang and Charley Lee and Leonard Tng and Serena Ge},
  title = {{DeepSWE}: Measuring Frontier Coding Agents on Original, Long-Horizon Engineering Tasks},
  year = {2026},
  eprint = {2607.07946},
  archivePrefix = {arXiv},
  primaryClass = {cs.SE},
  url = {https://arxiv.org/abs/2607.07946}
}
@misc{deepswev11,
  title = {DeepSWE v1.1: a cleaner, more reproducible benchmark for frontier coding agents},
  author = {Wenqi Huang and Peter Jiang},
  year = {2026},
  url = {https://github.com/datacurve-ai/deep-swe},
}
@misc{terminalbench21,
  author = {Mike A. Merrill and Alexander G. Shaw and Nicholas Carlini and others},
  title = {{Terminal-Bench}: Benchmarking Agents on Hard, Realistic Tasks in Command Line Interfaces},
  year = {2026},
  eprint = {2601.11868},
  archivePrefix = {arXiv},
  primaryClass = {cs.SE},
  url = {https://arxiv.org/abs/2601.11868}
}
@misc{terminalbench21release,
  author = {{The Terminal-Bench Team}},
  title = {{Terminal-Bench} 2.1},
  year = {2026},
  month = may,
  url = {https://www.tbench.ai/news/terminal-bench-2-1}
}
@misc{cwl,
  author = {Andrew Semenov and Svyatoslav Dorofeev},
  title = {Beyond Compaction: Structured Context Eviction for Long-Horizon Agents},
  year = {2026},
  eprint = {2606.11213},
  archivePrefix = {arXiv},
  primaryClass = {cs.CL},
  url = {https://arxiv.org/abs/2606.11213}
}
@misc{selfgc,
  author = {Xubin Hao and Hongjin Meng and Xin Yin and Jiawei Zhu and Chenpeng Cao},
  title = {Self-{GC}: Self-Governing Context for Long-Horizon {LLM} Agents},
  year = {2026},
  eprint = {2607.00692},
  archivePrefix = {arXiv},
  primaryClass = {cs.AI},
  url = {https://arxiv.org/abs/2607.00692}
}
@misc{vista,
  author = {Binyan Xu and Haitao Li and Kehuan Zhang},
  title = {{LLM} Agents Are Latent Context Managers: Eliciting Self-Managed Context via State Proprioception},
  year = {2026},
  eprint = {2606.30005},
  archivePrefix = {arXiv},
  primaryClass = {cs.CL},
  url = {https://arxiv.org/abs/2606.30005}
}
@article{lostmiddle,
  author = {Liu, Nelson F. and Lin, Kevin and Hewitt, John and Paranjape, Ashwin and Bevilacqua, Michele and Petroni, Fabio and Liang, Percy},
  title = {Lost in the Middle: How Language Models Use Long Contexts},
  journal = {Transactions of the Association for Computational Linguistics},
  volume = {12},
  pages = {157--173},
  month = feb,
  year = {2024},
  doi = {10.1162/tacl_a_00638},
  url = {https://aclanthology.org/2024.tacl-1.9/}
}
@inproceedings{toolsandbox,
  author = {Jiarui Lu and Thomas Holleis and Yizhe Zhang and Bernhard Aumayer and Feng Nan and Haoping Bai and Shuang Ma and Shen Ma and Mengyu Li and Guoli Yin and Zirui Wang and Ruoming Pang},
  title = {{ToolSandbox}: A Stateful, Conversational, Interactive Evaluation Benchmark for {LLM} Tool Use Capabilities},
  booktitle = {Findings of the Association for Computational Linguistics: NAACL 2025},
  pages = {1160--1183},
  year = {2025},
  doi = {10.18653/v1/2025.findings-naacl.65},
  url = {https://aclanthology.org/2025.findings-naacl.65/}
}
@inproceedings{toolemu,
  author = {Yangjun Ruan and Honghua Dong and Andrew Wang and Silviu Pitis and Yongchao Zhou and Jimmy Ba and Yann Dubois and Chris Maddison and Tatsunori Hashimoto},
  title = {Identifying the Risks of {LM} Agents with an {LM}-Emulated Sandbox},
  booktitle = {International Conference on Learning Representations},
  year = {2024},
  url = {https://proceedings.iclr.cc/paper_files/paper/2024/hash/7274ed909a312d4d869cc328ad1c5f04-Abstract-Conference.html}
}
@inproceedings{selectivecontext,
  author = {Li, Yucheng and Dong, Bo and Guerin, Frank and Lin, Chenghua},
  title = {Compressing Context to Enhance Inference Efficiency of Large Language Models},
  booktitle = {Proceedings of the 2023 Conference on Empirical Methods in Natural Language Processing},
  pages = {6342--6353},
  month = dec,
  year = {2023},
  doi = {10.18653/v1/2023.emnlp-main.391},
  url = {https://aclanthology.org/2023.emnlp-main.391/}
}
@article{sokolova2009performance,
  author = {Sokolova, Marina and Lapalme, Guy},
  title = {A Systematic Analysis of Performance Measures for Classification Tasks},
  journal = {Information Processing \& Management},
  volume = {45},
  number = {4},
  pages = {427--437},
  month = jul,
  year = {2009},
  doi = {10.1016/j.ipm.2009.03.002},
  url = {https://doi.org/10.1016/j.ipm.2009.03.002}
}
\end{filecontents*}
% Comment out the next line to use the anonymous AAMAS submission branch.
\def\ARCPreprintVersion{}
%%%%%%%%%%%%%%%%%%%%%%%%%%%%%%%%%%%%%%%%%%%%%%%%%%%%%%%%%%%%%%%%%%%%%%%%
%%% LaTeX Template for AAMAS-2027 (based on sigconf.tex)
%%% Prepared by the AAMAS-2027 Publication Chairs based on the version from AAMAS-2026.
%%%%%%%%%%%%%%%%%%%%%%%%%%%%%%%%%%%%%%%%%%%%%%%%%%%%%%%%%%%%%%%%%%%%%%%%
%%% Start your document with the \documentclass command.
%%% == IMPORTANT ==
%%% Use the first variant below to anonymize your submission (no author information shown).
%%% Use the second variant below for the final paper (including author information).
%%% For further information on anonymity and double-blind reviewing,
%%% please consult the call for paper information
%%% https://warwick.ac.uk/fac/sci/dcs/aamas2027/calls/
%%%% The default entry point builds the anonymized submission. Thin wrappers
%%%% select the authored conference and conference-neutral preprint versions.
\ifdefined\ARCPreprintVersion
  \documentclass[sigconf,nonacm]{aamas}
\else\ifdefined\ARCFormalVersion
  \documentclass[sigconf]{aamas}
\else
  \documentclass[sigconf,anonymous]{aamas}
\fi\fi
%%% Load required packages here (note that many are included already).
\usepackage{balance} % for balancing columns on the final page
\usepackage{amsmath}
\usepackage{amssymb}
\usepackage{mathtools}
\usepackage{booktabs}
\usepackage{multirow}
\usepackage{tabularx}
\usepackage{array}
\usepackage{xcolor}
\usepackage{tikz}
\usepackage{enumitem}
\usepackage{microtype}
\usepackage{graphicx}
\usepackage{float}
\setlength{\emergencystretch}{3em}
%%%%%%%%%%%%%%%%%%%%%%%%%%%%%%%%%%%%%%%%%%%%%%%%%%%%%%%%%%%%%%%%%%%%%%%%
%%% AAMAS-2027 metadata and copyright block (conference versions only).
\ifdefined\ARCPreprintVersion
  \setcopyright{none}
  \settopmatter{printacmref=false}
\else
\makeatletter
\gdef\@copyrightpermission{
  \begin{minipage}{0.2\columnwidth}
   \href{https://creativecommons.org/licenses/by/4.0/}{\includegraphics[width=0.90\textwidth]{by}}
  \end{minipage}\hfill
  \begin{minipage}{0.8\columnwidth}
   \href{https://creativecommons.org/licenses/by/4.0/}{This work is licensed under a Creative Commons Attribution International 4.0 License.}
  \end{minipage}
  \vspace{5pt}
}
\makeatother
\setcopyright{ifaamas}
\acmConference[AAMAS 2027]{Proc.\@ of the 26th International Conference
on Autonomous Agents and Multiagent Systems (AAMAS 2027)}{3--7 May 2027}
{Hanoi, Vietnam}{M.~Baldoni, F.~Fang, W.~Yeoh, N.~Yorke-Smith (eds.)}
\copyrightyear{2027}
\acmYear{2027}
\acmDOI{}
%\acmPrice{}
\acmISBN{}
%%%%%%%%%%%%%%%%%%%%%%%%%%%%%%%%%%%%%%%%%%%%%%%%%%%%%%%%%%%%%%%%%%%%%%%%
%%% == IMPORTANT ==
%%% Use this command to specify your OpenReview submission number.
%%% In anonymous mode, it will be printed on the first page.
% \acmSubmissionID{<OpenReview submission id>}
%%% Use this command to specify the track.
\submissionType{Research Paper Track}
\fi
%%% Use this command to specify the title of your paper.
\begin{document}
\title{ARC: A Contract-Governed Runtime for Semantic Continuity in Long-Horizon Agent Execution}
%%% Provide names, affiliations, and email addresses for all authors.
\ifdefined\ARCFormalVersion
  \settopmatter{authorsperrow=4}
\else\ifdefined\ARCPreprintVersion
  \settopmatter{authorsperrow=4}
\else
  \settopmatter{authorsperrow=2}
\fi\fi
% Preserve the exact mailto target while allowing one long address to wrap in
% the fixed-width four-author layout.
\makeatletter
\newcommand{\ARCbreakableemail}[3]{%
  \if@ACM@anonymous\else
    \g@addto@macro\addresses{\ARC@typesetbreakableemail{#1}{#2}{#3}}%
  \fi}
\newcommand{\ARC@typesetbreakableemail}[3]{%
  \ifx\@currentaffiliation\@empty
    \gdef\@currentaffiliation{\href{mailto:#1}{#2@\allowbreak #3}}%
  \else
    \g@addto@macro\@currentaffiliation{\par\href{mailto:#1}{#2@\allowbreak #3}}%
  \fi}
\makeatother
\author{Yicheng Du}
\affiliation{
  \institution{University of Science and Technology of China}
  \city{Hefei}
  \country{China}}
\email{dycalo@mail.ustc.edu.cn}
\author{Chenwei Wu}
\affiliation{
  \institution{Huawei Technologies Ltd.}
  \city{Dongguan}
  \country{China}}
\email{wuchenwei@huawei.com}
\author{Zuyu Zhao}
\affiliation{
  \institution{Huawei Technologies Ltd.}
  \city{Dongguan}
  \country{China}}
\email{zhaozuyu1@huawei.com}
\author{Weijian Sun}
\affiliation{
  \institution{Huawei Technologies Ltd.}
  \city{Shanghai}
  \country{China}}
\email{sunweijian@huawei.com}
\author{Dandan Miao}
\affiliation{
  \institution{Huawei Technologies Ltd.}
  \city{Shanghai}
  \country{China}}
\email{miaodandan@huawei.com}
\author{Mingyue Li}
\affiliation{
  \institution{Huawei Technologies Ltd.}
  \city{Beijing}
  \country{China}}
\ARCbreakableemail{limingyue54735@huawei.com}{limingyue54735}{huawei.com}
\author{Junjie Hu}
\affiliation{
  \institution{Southeast University}
  \city{Nanjing}
  \country{China}}
\email{jjhu\_seu@seu.edu.cn}
\author{Tong Xu\texorpdfstring{\textsuperscript{*}}{*}}
\affiliation{
  \institution{University of Science and Technology of China}
  \city{Hefei}
  \country{China}}
\email{tongxu@ustc.edu.cn}
%%% Use this environment to specify a short abstract for your paper.
\begin{abstract}
Long-horizon autonomous agents struggle to maintain sound context as intermediate artifacts accumulate over extended trajectories.
Prevailing context-management approaches rely on heuristic retrieval, summarization, or compression; however, they optimize for
textual relevance and context availability without formally specifying evidence obligations required to authorize a decision,
nor do they ensure that this evidence remains valid at commit time.
This disconnect between reasoning-time context and execution-time state breaks semantic continuity, causing
agents to act on incomplete, ungrounded, or stale information.
To address this challenge, we propose the Agent Runtime for Semantic Continuity (ARC),
a contract-governed execution framework that systematically bridges an agent’s reasoning context with dynamic execution environments.
Driven by declarative Domain Contracts, ARC coordinates context management across the decision lifecycle.
During reasoning, it synthesizes verifiable, budget-bounded context views that satisfy explicit evidence obligations and
reconstructs missing prerequisites lost to memory eviction. Before executing side effects,
ARC enforces commit-time validation by checking action dependencies against the live execution state,
aborting stale decisions caused by environmental drift.
Experiments across synthetic dataflows, IP network operations, and complex software engineering benchmarks
show that ARC maintains evidence completeness under context pressure, reduces unsafe commits,
and achieves high task completion rates while reducing peak context size and token overhead.
\end{abstract}
\ccsdesc[500]{Computing methodologies~Intelligent agents}
\ccsdesc[300]{Software and its engineering~Software architectures}
\keywords{LLM, Agent Runtimes, Context Management, Long-Horizon Tasks}
%%%%%%%%%%%%%%%%%%%%%%%%%%%%%%%%%%%%%%%%%%%%%%%%%%%%%%%%%%%%%%%%%%%%%%%%
%%% Include any author-defined commands here.
\newcommand{\method}{\textsc{ARC}}
\newcommand{\methodCRI}{\textsc{ARC-CRI}}
\newcommand{\Eligible}{\operatorname{Eligible}}
\newcommand{\Dep}{\operatorname{Dep}}
\newcommand{\Complete}{\operatorname{Complete}}
\newcommand{\Cost}{\operatorname{Cost}}
\newcommand{\Valid}{\operatorname{Valid}}
\newcommand{\Sat}{\operatorname{Sat}}
\newcommand{\Authorize}{\operatorname{Authorize}}
%%%%%%%%%%%%%%%%%%%%%%%%%%%%%%%%%%%%%%%%%%%%%%%%%%%%%%%%%%%%%%%%%%%%%%%%
%%% The following commands remove the headers in your paper. For final
%%% papers, these will be inserted during the pagination process.
\ifdefined\ARCPreprintVersion
  \pagestyle{plain}
\else
  \pagestyle{fancy}
  \fancyhead{}
\fi
%%% The next command prints the information defined in the preamble.
\maketitle
\ifdefined\ARCPreprintVersion
  \thispagestyle{plain}
  \hypersetup{pdfcreator={LaTeX}}
\fi
%%%%%%%%%%%%%%%%%%%%%%%%%%%%%%%%%%%%%%%%%%%%%%%%%%%%%%%%%%%%%%%%%%%%%%%%
\begin{figure}[t]
\centering
\includegraphics[width=\columnwidth]{figures/arc_graphical_abstract.png}
\Description{A comparison between transcript-centric context management and
ARC. The first path leaves the model responsible for history management and
does not bind evidence to side-effect authorization. The ARC path separates
prospective context requirements from runtime resolution while one Domain
Contract governs evidence admission and commit authorization.}
\caption{\textbf{(a)} Transcript-centric methods leave
the model entangled with history management and do not bind admitted
evidence to effect authorization. \textbf{(b)} ARC separates
\emph{what must be known next} from runtime decisions, while the Domain Contract spans context
admission and commit.}
\label{fig:graphical-abstract}
\end{figure}
\section{Introduction}
Long-horizon language-model agents increasingly interact with complex, evolving environments \cite{react,taubench}.
Their decisions may depend on artifacts produced many steps earlier, while only a small fraction of the execution history fits within
a bounded reasoning context. Consequently, critical evidence may be omitted from the model-visible view,
and facts observed during reasoning may become stale before the resulting side effect is committed.
Stateful tool-use benchmarks further expose implicit dependencies across tools and turns, with state-dependency tasks remaining challenging even for strong models \cite{toolsandbox}.
Even when relevant evidence remains within the nominal context window, models may use it unreliably as its position changes in a long input \cite{lostmiddle}.
Figure~\ref{fig:graphical-abstract} summarizes the resulting control-boundary gap and ARC's proposed separation of responsibilities.
Prior work addresses the challenges of bounded contexts and dynamic environments at disparate system boundaries.
Context systems manage interaction histories or source documents through paging, compression, hierarchy, and linked notes \cite{memgpt,readagent,hiagent,amem};
optimize compression or folding \cite{acon,agentfold}; maintain dependency graphs \cite{contextweaver} or typed dependency-linked episodes \cite{cwl};
or use learned or model-directed lifecycle and garbage-collection operations \cite{mem1,selfgc,vista}.
Runtime enforcement applies declarative policies \cite{agentspec,forge}, knowledge-enabled action guards \cite{guardagent}, or evidence-dependency checks \cite{groundedcontinuation};
transactional runtimes provide compensation \cite{sagallm,atomix}, read-set validation, or commit-time authority checks \cite{sbus,committimeauth}.
However, these lines of work do not expose a shared, decision-specific semantics that
persists from context admission to effect authorization.
This leaves a fundamental disconnect between what an agent is allowed to reason from and what it is allowed to commit,
which typically manifests in two primary failure modes. First, in reasoning-time witness omission,
existing context management policies optimize relevance or compression independently of governed decision requirements,
frequently omitting witnesses required by the active decision and precipitating hallucinations or downstream execution failures.
Second, in commit-time dependency invalidation, current execution gates validate live preconditions or action-side read sets without
binding their scope to the evidence admitted during reasoning. When the environment evolves between the reasoning snapshot and
the execution state, the decision's premises become obsolete, yet the runtime lacks the protocol-level provenance required to
detect that the action is no longer authorized.
To bridge this gap, we propose \method{}, a runtime architecture that establishes contract-governed semantic continuity across
the agent lifecycle by unifying context compilation and state validation via declarative Domain Contracts.
ARC makes the Agent--Runtime turn boundary explicit through a typed Context Requirement Interface (CRI), where each step returns an action alongside a prospective,
typed requirement set for the next invocation. After a proposal commits, the trusted requirement-state normalizer combines its activated declaration with requirements inferred from
managed provenance, observed access traces, and active global requirements. The Context Broker then resolves this normalized plan to candidate evidence under the available budget,
while the Domain Contract prevents the model from weakening required obligations. During reasoning,
\method{} compiles a budget-bounded view and emits a certificate attesting that every contract-specified obligation is represented by
an admissible witness, automatically triggering same-snapshot reconstruction if compression damages required evidence.
When the agent emits an action, the runtime projects the managed-resource references of the certified view together with
contract-defined action dependencies into a sealed validation scope, revalidating every dependency within this scope and
evaluating live preconditions before authorizing side effects at execution time.
% We evaluate \method{} across diverse domains spanning proprietary benchmarks in
% compositional dataflow inversion and IP network operations, alongside standard public benchmarks including
% DeepSWE and Terminal-Bench. By systematically varying context pressure, compression damage, and dynamic state mutations,
% \method{} improves governed task completion, blocks unsafe commits without penalizing valid actions,
% and reduces token and peak-context overhead on public benchmarks, demonstrating both the benefits and operational limits.
The primary contributions of this work are:
\begin{enumerate}
    \item We formulate \textbf{contract-governed semantic continuity} and \textbf{prospective requirement declaration}. The runtime normalizes model-emitted, inferred, observed, and global requirement signals, after which the Context Broker resolves a just-in-time context closure under explicit contract and budget constraints.
    \item We introduce \textbf{contract-admissible evidence views}. A budget-aware compiler constructs decision-specific views, which an independent verifier certifies; failed verification triggers same-snapshot reconstruction to recover from compression-induced witness omissions.
    \item We introduce a \textbf{certificate-bound action protocol}. It projects certified reasoning views and action semantics into a sealed dependency set, revalidating live execution state to intercept actions predicated on stale evidence.
\end{enumerate}
\section{Related Work}
\subsection{Agent Loops and Context Management}
ReAct interleaves reasoning and action traces, exposing tool observations directly to subsequent reasoning turns~\cite{react}.
To scale agents over long horizons, subsequent architectures introduced paging, virtual memory, hierarchical subgoals,
episodic compression, and learned context retention policies~\cite{memgpt,readagent,hiagent,amem,acon,agentfold,mem1}.
Recent systems depart from flat transcripts to expose structured context lifecycles.
CWL partitions execution into typed exploration and action episodes, allowing downstream actions to depend on closed exploration subgraphs
while driving graduated, budget-triggered eviction~\cite{cwl}.
ContextWeaver retrospectively extracts completed Thought--Action--Observation nodes, infers parent dependencies via an auxiliary model,
and retrieves validated ancestry during context construction~\cite{contextweaver}.
Similarly, Self-GC employs a side-channel planner to fold, mask, or prune indexed context objects through
a plan--rehearse--commit lifecycle~\cite{selfgc}, whereas VISTA exposes token budgets and archive dashboards to
facilitate highly flexible and imperative context block recovery operations~\cite{vista}.
ARC differs from these approaches in its prospective, runtime-mediated control boundary.
Rather than relying on retrospective graph extraction or imperative model-directed memory edits, ARC introduces
a declarative Context Requirement Interface.
The task model emits typed requirements for its \emph{next} invocation alongside its current action payload.
After successful commit, the trusted requirement-state normalizer combines the activated declaration with provenance-inferred, observed, and global requirements. The Context Broker then maps the normalized plan to candidate records,
resolving sources, fidelity, and budget allocations.
Furthermore, the resulting view is certified against declarative Domain Contracts and bound to commit-time authorization,
establishing a robust semantic continuity across the agent lifecycle that standalone context managers inherently lack.
\subsection{Safe and Transactional Execution}
A parallel line of work enforces safety, alignment, and correctness during agent execution.
Runtime enforcement frameworks such as AgentSpec, Forge, GuardAgent, and grounded-continuation systems monitor agent behaviors against
formal specifications, temporal rules, domain knowledge, or dynamic support graphs~\cite{agentspec,forge,guardagent,groundedcontinuation}.
These mechanisms gate actions or conversational continuations before release.
Transactional runtimes complement them: SagaLLM and Atomix provide settlement, compensation, and rollback~\cite{sagallm,atomix};
S-Bus reconstructs commit-time read sets~\cite{sbus}; CommitGuard revalidates effect-bound authority at durability~\cite{committimeauth}.
Despite these guarantees, existing methods operate over an incomplete verification boundary.
Policy monitors evaluate isolated action payloads against static rules,
remaining agnostic to whether the evidence underlying the agent's reasoning has degraded.
Conversely, transactional runtimes validate live preconditions and action-side read sets at execution time,
but do not bind these checks to the historical evidence admitted into the prompt.
To address this gap, our formulation draws on classical design-by-contract and
runtime-monitoring traditions~\cite{meyer1992contract,schneider2000enforcement,havelund2002monitors},
while integrating optimistic concurrency, multiversion control,
and data provenance for evidence tracking~\cite{kung1981occ,bernstein1983mvcc,buneman2001provenance}.
Building on these foundations, ARC bridges this divide by
making Domain Contracts the authoritative semantic bridge across the agent lifecycle.
The same contract that certifies evidence admissibility during reasoning also dictates
the mandatory baseline of action dependencies that must be revalidated at execution time.
This unified design ensures that commit authorization verifies not only live state validity,
but also the continuous freshness of the premises from which the decision was derived.
\section{Problem Formulation}
\label{sec:problem}
This section formulates the single-agent decision process over a bounded reasoning context and an evolving state. The formal objective is to enforce a safety property: the runtime must either construct an admissible reasoning view and authorize a proposal against the live state, or reject the proposal without application. All guarantees are defined relative to an instantiated Domain Contract and the runtime's managed resources. Table~\ref{tab:notation} summarizes the core notation.
\begin{table*}[htbp]
    \centering
    \small
    \setlength{\tabcolsep}{4pt}
    \caption{Core notation used in the formulation and method.}
    \label{tab:notation}
    \begin{tabular}{@{}p{0.15\textwidth}p{0.31\textwidth}
                        p{0.15\textwidth}p{0.31\textwidth}@{}}
    \toprule
    \textbf{Symbol} & \textbf{Meaning}
    & \textbf{Symbol} & \textbf{Meaning} \\
    \midrule
    $d,\mathcal C_d$ &
    Decision type and its active Domain Contract.
    & $\mathcal X,\nu_{\sigma}(x)$ &
    Managed resources and the version of resource $x$ at snapshot $\sigma$.
    \\
    $\sigma_s,\sigma_e$ &
    Reasoning snapshot and commit-time snapshot for one decision.
    & $\mathcal H_s,V_s,B$ &
    Available records, certified evidence view, and rendering budget.
    \\
    $\widehat Q_s,Q^{\mathrm{decl}}$ &
    Normalized requirements for an invocation and its pending next-turn declaration.
    & $O_d,\Gamma_d,\Delta_d,P_d$ &
    Evidence obligations, view-admission rule, action-dependency rule, and live precondition.
    \\
    $W_o,\pi_s$ &
    Witness records for obligation $o$ and the resulting view certificate.
    & $\alpha,R_{\mathrm{add}}$ &
    Action payload and model-declared additional guarded resources.
    \\
    $a,D(a),\bot$ &
    One-shot sealed proposal, its commit-validation dependencies, and rejection.
    & $C_s$ &
    Exact model input at the reasoning snapshot.
    \\
    $\operatorname{Render},\operatorname{Cost},\operatorname{Refs}$ &
    Canonical rendering, budget cost, and versioned references carried by a view.
    & $\operatorname{Adm},\operatorname{Verify},$
      $\operatorname{Bound},\Authorize$ &
    View admissibility, certificate verification, proposal binding, and authorization.
    \\
    \bottomrule
    \end{tabular}
\end{table*}
\subsection{Setting and Governed-Decision Task}
Let $\mathcal X$ be a set of managed resources. A snapshot $\sigma$ assigns each resource $x \in \mathcal X$ a version $\nu_{\sigma}(x)$. Let $\sigma_s$ and $\sigma_e$ denote the reasoning and commit-time snapshots for a given decision, where $\nu_s(x)=\nu_{\sigma_s}(x)$ and $\nu_e(x)=\nu_{\sigma_e}(x)$. At $\sigma_s$, the runtime holds an immutable record set $\mathcal H_s$. Each record contains a stable identifier, renderable content, provenance, and transitive resource--version dependencies. $\operatorname{Refs}(V_s)$ denotes the union of these dependencies for records within a view $V_s\subseteq\mathcal H_s$.
The model-visible context exclusively consists of static instructions and the rendered view of managed records:
\begin{equation}
C_s=C_s^{\mathrm{static}}\oplus\operatorname{Render}(V_s).
\label{eq:certified-context}
\end{equation}
Here, $\operatorname{Render}$ is deterministic and $\oplus$ denotes ordered concatenation. Consequently, any managed-state premise exposed to the model must be included in $V_s$ and carry versioned references.
For a decision type $d$, the active Domain Contract is formalized as:
\begin{equation}
\mathcal C_d=
\left(O_d,\Gamma_d,\Delta_d,P_d\right).
\label{eq:domain-contract}
\end{equation}
Each $o\in O_d$ is a runtime-checkable predicate over a witness bundle and a snapshot. The view-admission rule $\Gamma_d(V,\widehat Q;\sigma)$ returns $1$ if the records possess valid provenance and authority, and satisfy the required items in $\widehat Q$ at the specified representation. $\Delta_d(\alpha,\sigma)$ defines the versioned dependencies of action payload $\alpha$, and $P_d(\alpha,\sigma)$ specifies its live precondition. These rules are explicitly defined by the domain specification.
Given $(d,\mathcal C_d,\mathcal H_s,\widehat Q_s,\sigma_s,B)$, the runtime coordinates at most one model invocation and one side effect via the following protocol:
\begin{equation}
\begin{aligned}
(\mathcal H_s,\widehat Q_s,\sigma_s,B)
&\longrightarrow(V_s,\pi_s)\\
&\longrightarrow(\alpha,R_{\mathrm{add}},Q^{\mathrm{decl}})
\longrightarrow(a,D(a))\\
&\longrightarrow\{\mathsf{commit},\bot\}.
\end{aligned}
\label{eq:governed-task}
\end{equation}
The protocol ensures that the model is invoked exclusively on a contract-admissible certified view, and that a proposal is committed only if it binds to that invocation, its dependencies remain fresh, and its live precondition is satisfied. The protocol is fail-closed; view construction, output validation, or authorization may yield $\bot$. A pending declaration $Q^{\mathrm{decl}}$ remains inactive until its associated proposal successfully commits.
\subsection{Contract-Admissible Reasoning Evidence}
A view is contract-admissible if it is drawn from the available snapshot, satisfies the rendering budget $B$, meets the view-level admission rule, and contains valid witness bundles $W_o$ for all mandatory obligations:
\begin{equation}
\begin{aligned}
&\operatorname{Adm}_{\mathcal C_d}
(V_s,\widehat Q_s,\sigma_s;B)=1\\
&\quad\Longleftrightarrow
V_s\subseteq\mathcal H_s\\
&{}\land\left[\forall(x,\bar\nu)\in\operatorname{Refs}(V_s):
\nu_s(x)=\bar\nu\right]\\
&{}\land \operatorname{Cost}(\operatorname{Render}(V_s))\le B\\
&{}\land \Gamma_d(V_s,\widehat Q_s;\sigma_s)=1\\
&{}\land \forall o\in O_d\;\exists W_o\subseteq V_s:
o(W_o;\sigma_s)=1 .
\end{aligned}
\label{eq:admissible-view}
\end{equation}
The $\operatorname{Cost}$ function maps rendering outputs to domain-specific limits (e.g., token counts). Required items in $\widehat Q_s$ function as hard constraints in $\Gamma_d$, while optional items influence selection without affecting admissibility. Equation~\eqref{eq:admissible-view} strictly enforces snapshot consistency.
For an admitted view, the runtime issues a protected certificate $\pi_s$. This certificate binds the contract version, invocation, snapshot, normalized requirements, budget, witness choices, record manifest, and exact rendering. The model is constrained to the input in Equation~\eqref{eq:certified-context} and cannot modify $\pi_s$. An independent verifier enforces the following soundness condition:
\begin{equation}
\operatorname{Verify}(\pi_s,\mathcal C_d)=1
\Longrightarrow
\operatorname{Adm}_{\mathcal C_d}
(V_s,\widehat Q_s,\sigma_s;B)=1.
\label{eq:admission-soundness}
\end{equation}
This condition guarantees contract-relative evidence coverage for a specific invocation; it evaluates formal admissibility independent of reasoning correctness or unspecified external dependencies.
\subsection{Evidence-Bound Commit Authorization}
\label{sec:evidence-bound-actions}
The model outputs an action payload $\alpha$, a set $R_{\mathrm{add}}\subseteq\mathcal X$ of supplementary managed resources whose current versions must be guarded, and a declaration $Q^{\mathrm{decl}}$. The runtime encapsulates this tuple into a one-shot proposal $a$ bound to $\pi_s$, sealing the validation scope $D(a)$ as follows:
\begin{equation}
D(a)=\operatorname{Refs}(V_s)
 \cup\Delta_d(\alpha,\sigma_s)
 \cup\{(x,\nu_s(x))\mid x\in R_{\mathrm{add}}\}.
\label{eq:action-dependencies}
\end{equation}
Thus, $D(a)$ structurally lower-bounds dependencies using the certified evidence and contract rules, while model-declared guarded resources may expand, but never reduce, this set.
The binding predicate $\operatorname{Bound}_{\mathcal C_d}(a,\pi_s)$ verifies that $a$ is schema-valid, derives from the invocation tied to $\pi_s$, aligns with the exact runtime-computed dependencies in Equation~\eqref{eq:action-dependencies}, and remains unspent. At the side-effect linearization point, authorization requires:
\begin{equation}
\begin{aligned}
&\Authorize_{\mathcal C_d}(a,\pi_s,\sigma_e)=1
\quad\Longleftrightarrow\quad
\operatorname{Verify}(\pi_s,\mathcal C_d)=1\\
&\qquad{}\land
\operatorname{Bound}_{\mathcal C_d}(a,\pi_s)\\
&\qquad{}\land
\left[\forall(x,\bar\nu)\in D(a):\nu_e(x)=\bar\nu\right]\\
&\qquad{}\land P_d(\alpha,\sigma_e)=1 .
\end{aligned}
\label{eq:commit-authorization}
\end{equation}
The Commit Gate performs authorization and the resulting terminal transition atomically. If $\Authorize_{\mathcal C_d}=1$ and the managed application succeeds, it applies the side effect, consumes the proposal, and activates the pending $Q^{\mathrm{decl}}$ at the same linearization point; otherwise, it rejects the proposal without a managed effect and discards the pending declaration. Atomic proposal consumption structurally prevents replay.
\subsection{Safety Objective and Scope}
\noindent\textbf{Proposition 1 (Contract-governed semantic continuity).}
Assume verifier soundness (Equation~\eqref{eq:admission-soundness}), protected runtime certificates and sealing, consistent resource versioning, type-correct contract rules, and atomic authorization--application--consumption--activation. Every successfully committed proposal originates from a contract-admissible view at $\sigma_s$, is strictly bound to its invocation, observes invariant versions for all $D(a)$ dependencies at $\sigma_e$, and satisfies $P_d(\alpha,\sigma_e)$.
\emph{Proof sketch.} A successful commit strictly entails Equation~\eqref{eq:commit-authorization}. Verification guarantees view admissibility via Equation~\eqref{eq:admission-soundness}. Binding ensures the association with the invocation and its sealed dependencies. The remaining predicates establish freshness and the precondition. Atomic consumption prevents repeated applications.
This guarantee is bounded by the defined managed state. Assuming complete specification of $\Delta_d$ for relevant operations, the protocol enforces semantic continuity for managed resources. Contract predicates are executable runtime constraints derived from domain specifications; they operate independently of benchmark-oracle labels ($E_i$ and $A_i$) utilized in Section~\ref{sec:setup}. Atomicity is structurally required; without it, the property reduces to standard pre-dispatch validation.
\section{Method}
\label{sec:method}
\method{} implements Equation~\eqref{eq:governed-task} through prospective requirement synthesis, contract-guided view certification, and certificate-bound authorization. Figure~\ref{fig:arc-framework} outlines the execution pipeline. Structurally, the model proposes semantic requirements and additional guarded resources, while the runtime exclusively handles requirement normalization, evidence selection, certificate issuance, dependency sealing, and authorization.
The trusted computing base is restricted to the contract registry and requirement-state normalizer, versioned record store, canonical renderer, independent verifier, runtime sealer, Commit Gate, and atomic conditional-execution backend. The Context Broker, model, compressor, and view compiler are untrusted. ARC constrains the evidence used for proposal generation and the state against which the proposal is executed, without verifying the internal reasoning correctness of the model.
\begin{figure*}[htbp]
\centering
\includegraphics[width=0.8\textwidth]{figures/arc_framework_detailed_white.png}
\Description{A white-background ARC architecture with four parallel stages. Requirement signals are normalized and resolved to candidate evidence; an independently verified view enters the Agent while its protected certificate bypasses the Agent and enters the Runtime Sealer; the sealed proposal then reaches commit-time authorization and atomic application. A versioned store lies below the stages, with separate next-turn, fresh-decision, and pre-invocation rebuild paths.}
\caption{\method{}'s contract-governed decision cycle. The Domain Contract supplies evidence-obligation, admission, dependency, and precondition rules. The trusted normalizer combines active declared, inferred, observed, and global requirements into $\widehat Q_s$, and the Context Broker resolves candidate records. A compiled view is independently verified into $(V_s,\pi_s)$: only $\operatorname{Render}(V_s)$ enters the Agent, while the protected certificate remains on the runtime path. The Agent's action, additional guarded resources, and pending next-turn declaration are bound by the Runtime Sealer into a one-shot proposal. Within the atomic commit transition, the Commit Gate checks certificate binding, dependency freshness, and the live precondition before applying and consuming an authorized proposal and activating its declaration. Rejection requires a fresh decision; the local recovery loop denotes pre-invocation rebuilding at the same snapshot.}
\label{fig:arc-framework}
\end{figure*}
\begin{figure}[tb]
\begin{center}
\begin{minipage}{0.9\columnwidth}
\hrule
\noindent\textbf{Algorithm 1: \method{} governed-decision cycle}
\footnotesize
\noindent\textbf{Input:} contract $\mathcal C_d$; snapshot
$(\sigma_s,\mathcal H_s)$; active declared, inferred, observed, and global
requirement state; budget $B$; Agent.\\
\noindent\textbf{Output:} committed observation with updated requirement
state, fail-closed result, or a fresh-snapshot re-decision.
\begin{tabularx}{\linewidth}{@{}r@{\hspace{0.45em}}X@{}}
1: & $\widehat Q_s\leftarrow\textsc{Normalize}$ (active declared, inferred,
     observed, and global requirements). \\
2: & Candidates $\leftarrow\textsc{ContextBroker}
     (\mathcal C_d,\widehat Q_s,\mathcal H_s,\sigma_s,B)$. \\
3: & At fixed $\sigma_s$, compile a candidate $V_s$ and, only after independent
     verification under Equation~\eqref{eq:admissible-view}, issue $\pi_s$. \\
4: & Upon materialization or verification failure, rematerialize at the same
     $\sigma_s$ and repeat Lines 2--3; return $\bot$ only when alternatives
     are exhausted. \\
5: & $(\alpha,R_{\mathrm{add}},Q^{\mathrm{decl}})\leftarrow
     \operatorname{Agent}(C_s^{\mathrm{static}}\oplus\operatorname{Render}(V_s))$. \\
6: & Validate the raw CRI tuple; upon failure, return $\bot$ without creating a
     proposal. Otherwise, seal it as a one-shot $a$ bound to $\pi_s$, its certified
     invocation, the immutable $D(a)$, and pending $Q^{\mathrm{decl}}$. \\
7: & At the side-effect linearization point, obtain live $\sigma_e$ and execute one
     atomic conditional transition under Equation~\eqref{eq:commit-authorization}:
     on success, apply $\alpha$, consume $a$, and activate $Q^{\mathrm{decl}}$;
     otherwise, reject $a$ without application and discard the declaration. \\
8: & After rejection, recompute the requirement plan at a fresh snapshot. \\
9: & After success, return the execution observation and updated requirement state. \\
\end{tabularx}
\hrule
\end{minipage}
\end{center}
\Description{Pseudocode for one ARC governed-decision cycle. The runtime normalizes requirements, compiles and independently verifies a snapshot-consistent view, invokes the agent, seals the proposed action with its certificate and dependencies, and atomically authorizes or rejects it at commit time. Pre-invocation failures rebuild at the same snapshot, whereas commit failures trigger a fresh-snapshot re-decision.}
\end{figure}
\subsection{Prospective Requirement Synthesis}
Let $t$ denote agent turns, while $s$ and $e$ indicate the reasoning and commit snapshots within a turn. Each agent response adheres to the Context Requirement Interface (CRI):
\begin{equation}
(\alpha_t,R_{\mathrm{add},t},Q^{\mathrm{decl}}_{t+1})
=\operatorname{Agent}(C_{s,t}).
\label{eq:cri-output}
\end{equation}
$C_{s,t}$ is instantiated via Equation~\eqref{eq:certified-context}. A declared requirement specifies a stable resource identifier, required/optional modality, full/summary/metadata representation level, and scope (one-shot or global). The field $Q^{\mathrm{decl}}$ is mandatory (though empty sets are valid), strictly decoupling semantic requirement formulation from the runtime's physical record selection.
The declaration remains pending with the proposal. A successful commit activates $Q^{\mathrm{decl}}_{t+1}$; a rejection discards both. Prior to the subsequent invocation, the runtime computes:
\begin{equation}
\widehat Q_{t+1}=\operatorname{Normalize}\!\left(
G_{t+1}\cup Q^{\mathrm{decl}}_{t+1}\cup Q^{\mathrm{inf}}_{t+1}
\cup\operatorname{Obs}(L_t)\right).
\label{eq:prospective-read-set}
\end{equation}
Here, $G_{t+1}$ contains active global requirements, $Q^{\mathrm{inf}}_{t+1}$ incorporates derived provenance and action semantics, and $\operatorname{Obs}(L_t)$ records runtime-level store and tool accesses. Normalization deterministically resolves conflicts: required overrides optional, finer representations override coarser ones, and global scope persists.
The Context Broker maps $\widehat Q_{t+1}$ to snapshot-consistent candidates under budget $B$. The resolver manages semantic localization, whereas the subsequent certification stage handles contract evidence formatting. Note that $Q^{\mathrm{decl}}_{t+1}$ defines future requirements, whereas $D(a)$ strictly dictates commit dependencies.
A generic baseline variant, \methodCRI{}, retains requirement resolution and budgeted materialization under a default eligibility policy. However, it omits the task-specific Domain Contract $\mathcal C_d$ and certificates, thereby not invoking Proposition 1.
\subsection{Contract-Guided View Synthesis and Certification}
\label{sec:view-construction}
Records in the append-only versioned store explicitly maintain identifiers, content, provenance, and transitive version dependencies. Derived artifacts (e.g., summaries or compressed inputs) are formalized as standard records, requiring verifiable provenance at $\sigma_s$. Eviction policies affect the model-visible cache, not authoritative snapshot records.
Given $(\mathcal C_d,\widehat Q_s,\sigma_s,B)$, the Context Broker's resolver produces candidate witness bundles for $O_d$ alongside supplemental records per $\Gamma_d$. The View Compiler searches for a $V_s$ satisfying Equation~\eqref{eq:admissible-view}. Among feasible candidates, selection maximizes contract-prioritized optional coverage, then minimizes cost, record count, and lexicographical identifier order. For additive cost functions, the compiler utilizes a deterministic greedy algorithm, optionally followed by exact mixed-integer programming (MIP) under identical constraints. Infeasibility results in $\bot$.
The compiler output is treated as untrusted. The independent verifier reconstructs the view using snapshot-bound handles, validates all obligations, coverage constraints, and budget limits, and deterministically recomputes the rendering. Upon success, a protected certificate $\pi_s$ is generated. The certificate binds cryptographic digests of the snapshot manifest and canonical rendering, while runtime-level protection ensures its integrity and authenticity. The model is subsequently invoked via Equation~\eqref{eq:certified-context}.
If a draft is invalidated prior to invocation (e.g., due to cache eviction or paging), the runtime rematerializes the requisite records from the authoritative store at $\sigma_s$ and repeats certification. A successful rebuild yields a new certificate. This recovery path occurs exclusively before model invocation; exhausting alternatives yields $\bot$.
\subsection{Certificate-Bound Authorization and Recovery}
\label{sec:commit-authorization}
Post-invocation, the Runtime Sealer validates the CRI tuple (Equation~\eqref{eq:cri-output}), resolves the versions of $R_{\mathrm{add}}$ at $\sigma_s$, and derives $D(a)$ via Equation~\eqref{eq:action-dependencies}. It then encapsulates the tuple as a one-shot proposal $a$, binding $\alpha$ and pending $Q^{\mathrm{decl}}$ to $\pi_s$, its invocation, and the immutable $D(a)$. The predicate $\operatorname{Bound}_{\mathcal C_d}(a,\pi_s)$ holds if and only if $a$ is schema-valid, unspent, bound to that exact certificate and invocation, and carries exactly the runtime-computed $D(a)$. This prevents downstream modification of certified dependencies.
For dynamic collections, $\Delta_d$ tracks the namespace or aggregate version to capture membership adjustments and in-place mutations. Additional guarded resources in $R_{\mathrm{add}}$ can only expand, and never reduce, the validation baseline $D(a)$.
At $\sigma_e$, the Commit Gate evaluates Equation~\eqref{eq:commit-authorization} and performs the terminal transition defined after that equation. Validation remains restricted to $D(a)$; changes to unreferenced managed resources do not intrinsically invalidate proposals, whereas excessive guarding may trigger conservative rejection. Any authorization or managed-application failure permanently rejects the proposal and discards its pending declaration.
If state drift induces rejection post-invocation, repairing historical views cannot restore authorization. The protocol instead initiates a fresh snapshot and recomputes the requirement plan without activating the rejected $Q^{\mathrm{decl}}_{t+1}$. This mechanism differs from the pre-invocation recovery (Section~\ref{sec:view-construction}): post-invocation rejection discards model outputs and advances to a fresh state, whereas pre-invocation recovery repairs evidence prior to generating an action.
\section{Experiment}
\label{sec:setup}
Our evaluation is designed to assess the effectiveness of contract-governed semantic continuity across the agent lifecycle.
Rather than merely measuring endpoint task success, we systematically investigate how \method{} manages context pressure,
detects compression-induced omissions and triggers fail-closed recovery, and blocks unsafe commits caused by environmental drift.
Furthermore, we evaluate whether the prospective-requirement loop transfers effectively to open-ended software and terminal tasks without
relying on task-specific hidden evidence.
\subsection{Implementation and Benchmarks}
We implement \method{} as a runtime middleware layer around \texttt{LangGraph} in Python.
The system relies on an append-only versioned record store, declarative contract checkers, deterministic view rendering, and
a Context Broker that resolves the versioned query plan between the action commit and the next model invocation.
Within each comparison, all methods share canonical task and domain semantics, actor, tool schemas, actor-context and cumulative-token ceilings, and sampling policy; semantics use each method's native carrier (e.g., contract, skill, or system prompt). Frozen method-native auxiliary-call ceilings may differ, but all auxiliary calls and tokens are charged to the originating method.
To evaluate \method{} across varying levels of environmental control and contract richness,
we select two governed project benchmarks and two public benchmarks:
\begin{enumerate}
    \item \textbf{Compositional Dataflow Inversion (CDI):} An 80-task benchmark where the agent infers and inverts a multi-stage transformation pipeline. A symbolic interpreter verifies transformation consistency and final-output equivalence.
    \item \textbf{IP Network Operations (IPO):} An 80-task suite for evidence-grounded physical-topology recovery, evaluated against a deterministic private topology oracle.
    \item \textbf{DeepSWE \cite{deepswebenchmark,deepswev11}:} We use the 113-task v1.1 snapshot. The agent operates without task-network access, and success is strictly determined by the official benchmark verifier.
    \item \textbf{Terminal-Bench 2.1 \cite{terminalbench21,terminalbench21release}:} We use the 89-task release snapshot, retaining official timeouts and resource constraints and evaluating solely with the benchmark's native verifier.
\end{enumerate}
Specifically, we utilize an internally deployed GLM-5.2-200K as the actor model for the governed evaluations (CDI and IPO),
and the provider alias deepseek-v4-flash, accessed on August 23--24, 2026, for the public benchmarks (DeepSWE and Terminal-Bench). Each method runs three trials per task; outcome metrics are trial-averaged within task and macro-averaged across tasks, with other denominators stated separately.
Project-specific Domain Contracts are authored from specifications prior to task assignment,
and their checkers operate independently of the hidden task oracles.
For the public benchmarks, we evaluate \methodCRI{}, the prospective context-requirement component without custom Domain Contracts or contract-governed safety claims. Their sole task-success authority remains the unmodified official verifier, ensuring a fair transfer evaluation.
\begin{table}[tb]
    \centering
    \footnotesize
    \setlength{\tabcolsep}{3pt}
    \caption{Benchmark assignments and primary authorities used in the evaluation.}
    \label{tab:benchmark-statistics}
    \begin{tabularx}{\columnwidth}{@{}l c >{\raggedright\arraybackslash}X >{\raggedright\arraybackslash}X@{}}
    \toprule
    \textbf{Benchmark} & \textbf{Tasks} & \textbf{Endpoint Authority} & \textbf{Execution Budget} \\
    \midrule
    CDI & 80 & Symbolic verifier & Main: $B/B_i^\star=2.0$; 10-point sweep \\
    IPO & 80 & Private topology \& oracle & 56k context tokens \\
    DeepSWE & 113 & Separate program verifier & 160 turns; 64k context \\
    Terminal-Bench & 89 & Benchmark native verifier & 160 turns; 64k context \\
    \bottomrule
    \end{tabularx}
\end{table}
\subsection{Baselines}
We compare \method{} against mechanism-matched baselines addressing both reasoning-side context management and
execution-side policy enforcement.
To isolate the impact of \method{}'s Commit Gate for authorization,
we evaluate alternative runtime policies using a standardized set of pre-generated action proposals.
The evaluated baselines include:
\begin{itemize}
    \item \textbf{ACON} \cite{acon}: Bounded rolling context-compression.
    \item \textbf{ContextWeaver} \cite{contextweaver}: Selective structured memory.
    \item \textbf{ReadAgent-P} \cite{readagent}: Episodic gist memories.
    \item \textbf{HiAgent} \cite{hiagent}: Subgoal-chunked working memory.
    \item \textbf{A-MEM} \cite{amem}: Dynamic structured memory notes.
    \item \textbf{Live-P Gate} \cite{agentspec}: Validation against runtime policies.
    \item \textbf{Global-V Gate}: Using global version constraints.
    \item \textbf{Read-Set Gate} \cite{sbus}: Based on the tracked action read-sets.
    \item \textbf{CommitGuard} \cite{committimeauth}: Commit-time authorization.
\end{itemize}
For fair comparison, all context policies receive the same task-visible evidence universe and provenance graph,
and all auxiliary operations (e.g., summarization, paging) are charged to the originating method. Baseline names denote mechanism-matched adapters in our common harness rather than unmodified reference systems.
Finally, we ablate view-to-action binding and commit revalidation using the same pre-generated proposals. The \emph{w/o View Bind} variant retains view certification and live action-side checks but excludes certified-view references from the sealed validation scope, whereas \emph{w/o Revalid.} retains proposal binding and the live precondition check but disables version-freshness validation over $D(a)$.
\subsection{Metrics}
Our evaluation utilizes a comprehensive set of metrics spanning end-to-end task performance, context semantic integrity, and commit safety.
For the CDI benchmark, the overarching metric for semantic continuity is \emph{Verifier-Backed Authorized Completion} (VBAC).
For an episode $i$, let $S_i \in \{0,1\}$ indicate endpoint task success verified by the oracle.
Let $E_i \in \{0,1\}$ indicate that every governed decision is a Complete Invocation: its oracle-defined obligations are assessed as complete by a method-blinded semantic judge, an actor invocation occurs, and its request lineage is valid.
Let $A_i \in \{0,1\}$ denote benchmark-oracle authorization validity, meaning that no action labelled invalid under the injected state mutations was committed. These are evaluation labels, not claims that the runtime knows the true dependency set or prevents model errors.
VBAC is defined as:
\begin{equation}
\operatorname{VBAC}
=
\frac{1}{N}
\sum_{i=1}^{N}
S_i E_i A_i .
\label{eq:vbac}
\end{equation}
To further dissect context integrity,
we measure \emph{Semantic Coverage} (the fraction of oracle-defined semantic obligations judged complete in the actor-visible view),
\emph{Complete View} (the percentage of governed decision points whose rendered context fulfills all oracle-labelled obligations),
and \emph{Complete Invocation} (the percentage that additionally proceed to an actor invocation with valid request lineage). Action correctness is evaluated separately by the independent symbolic or domain verifier and contributes through endpoint success $S_i$, not Complete Invocation.
Conversely, the \emph{Continuity-Violation Rate} (CVR) quantifies the percentage of committed governed actions
executing with missing dependencies.
To characterize commit-time authorization under simulated mutations,
we define the \emph{Unsafe Commit Rate} (UCR) and \emph{Valid-Action Preservation} (VAP):
\begin{equation}
\operatorname{UCR} = \frac{N(\mathrm{commit}\land\mathrm{invalid})}{N(\mathrm{invalid})}, \quad
\operatorname{VAP} = \frac{N(\mathrm{commit}\land\mathrm{valid})}{N(\mathrm{valid})}.
\label{eq:commit-metrics}
\end{equation}
We report UCR and VAP jointly to expose the safety--utility trade-off between rejecting invalid actions and preserving valid ones, consistent with tool-agent evaluations that assess safety alongside helpfulness \cite{toolemu}.
For IPO, we adapt VBAC to \emph{Verifier-Backed Workflow Completion} (VBWC) to measure overall workflow success, and \emph{Qualified-VBWC} (Q-VBWC) for runs that additionally pass strict topology auditing. For each trial, typed predicted and reference topology sets $P,G$ (canonical node identifiers and link keys) yield $F_1=2|P\cap G|/(|P|+|G|)$, following standard F-measure terminology \cite{sokolova2009performance}. We report their task-macro mean as \emph{Macro F1}; missing or invalid deliveries and zero denominators score zero by convention.
For the public benchmarks (DeepSWE and Terminal-Bench), we adopt their official metrics: \emph{Pass} denotes verifier-backed exact success, and \emph{Partial} is DeepSWE's diagnostic score for intermediate progress. \emph{Calls} reports mean provider invocations per assigned trial; public-benchmark \emph{Token} is the mean provider-reported input-plus-output usage per assigned trial; \emph{Tok./VBWC} normalizes total IPO usage by completed workflows; and \emph{Peak context} averages the within-trial maximum model-input footprint.
Reporting task outcomes together with token and peak-context footprints follows the efficiency--quality view used in context-compression evaluation \cite{selectivecontext}.
\section{Results}
\label{sec:results}
We first evaluate in-domain semantic continuity and commit-time robustness on CDI, then staged physical-topology recovery on the IPO benchmark, and finally transfer to public benchmarks.
\subsection{CDI Benchmark}
\begin{table}[tb]
    \centering
    \footnotesize
    \setlength{\tabcolsep}{3pt}
    \caption{CDI results (three trials/task). View and outcome percentages are task-macro; CVR is action-normalized; Calls is the per-trial mean.}
    \label{tab:e2e-results}
    \resizebox{\columnwidth}{!}{%
    \begin{tabular}{@{}l ccccccc @{}}
    \toprule
    \multirow{2}{*}{\textbf{Method}}
    & \textbf{Sem.}
    & \textbf{Comp.}
    & \textbf{Comp.}
    & \textbf{End.}
    & \multirow{2}{*}{\textbf{VBAC}}
    & \multirow{2}{*}{\textbf{CVR}}
    & \multirow{2}{*}{\textbf{Calls}} \\
    & \textbf{Cov.} & \textbf{View} & \textbf{Inv.} & \textbf{Succ.} & & & \\
    \midrule
    ACON          & 64.0 & 45.0 & 41.3 & 61.3 & 37.5 & 20.0 & 7.8 \\
    ContextWeaver & 82.8 & 72.5 & 60.0 & 57.5 & 56.3 & 10.0 & 13.2 \\
    ReadAgent-P   & 83.4 & 62.5 & 57.5 & 61.3 & 55.0 & 16.3 & 10.9 \\
    HiAgent       & 89.6 & 72.5 & 70.0 & 73.8 & 70.0 & 7.5  & 9.2 \\
    A-MEM         & 85.1 & 67.5 & 66.3 & 75.0 & 65.0 & 12.5 & 12.4 \\
    \midrule
    \textbf{\method{}} & \textbf{98.5} & \textbf{96.3} & \textbf{91.3} & \textbf{92.5} & \textbf{90.0} & \textbf{1.3} & \textbf{6.4} \\
    \bottomrule
    \end{tabular}%
    }
\end{table}
We first assess the core semantic continuity and execution safety on the CDI benchmark.
As shown in Table~\ref{tab:e2e-results}, \method{} achieves a VBAC of 90.0\%,
outperforming the strongest baseline HiAgent, while maintaining the lowest CVR at 1.3\%.
Taken together, these results show that \method{} jointly maintains evidence completeness and authorized task completion under the evaluated CDI workloads.
\begin{figure*}[htbp]
    \centering
    \includegraphics[width=\textwidth]{figures/rq1_context_pressure.pdf}
    \Description{Line plots compare ARC and five context-management baselines across ten measured normalized CDI budget levels. From left to right, the panels report Semantic Coverage, Complete Invocation, and Endpoint Success; each panel contains a vertical dashed line at the main-table budget of 2.0.}
    \caption{CDI context-pressure results across 10 measured normalized budgets. The panels report Semantic Coverage, Complete Invocation, and Endpoint Success. Budget is $B/B_i^\star$, where $B_i^\star$ is CDI's task-specific minimum; dashed lines mark the Table~\ref{tab:e2e-results} setting $B/B_i^\star=2.0$.}
    \label{fig:context-pressure}
\end{figure*}
Next, we investigate the systems' robustness against reasoning-time witness omission by varying the context budget constraint
(Figure~\ref{fig:context-pressure}).
Across the evaluated range, \method{} maintains consistently high Semantic Coverage, Complete Invocation, and Endpoint Success. Under tighter budgets, competing methods generally attain lower Complete Invocation and Endpoint Success and exhibit greater budget sensitivity. This robustness is consistent with \method{}'s fail-closed same-snapshot reconstruction policy when required evidence cannot be certified.
\begin{table}[tb]
    \centering
    \footnotesize
    \setlength{\tabcolsep}{4pt}
    \caption{CDI commit-time authorization safety and valid-action preservation; latency is in-process gate-only p50/p95.}
    \label{tab:commit-results}
    \resizebox{\columnwidth}{!}{%
    \begin{tabular}{@{}l ccccc @{}}
    \toprule
    \multirow{2}{*}{\textbf{Method}}
    & \multicolumn{2}{c}{\textbf{UCR} $\downarrow$}
    & \multicolumn{2}{c}{\textbf{VAP} $\uparrow$}
    & \textbf{Latency ($\mu$s)} \\
    \cmidrule(lr){2-3} \cmidrule(lr){4-5}
    & \textbf{Act-Dep} & \textbf{Precond}
    & \textbf{Irrel.} & \textbf{No Chg.}
    & \textbf{p50/p95} $\downarrow$ \\
    \midrule
    Live-P Gate             & 92.5 & 2.5 & 98.8 & 100  & 0.52/0.93 \\
    Global-V Gate           & 1.3  & 0   & 5.0  & 98.8 & 0.57/1.02 \\
    Read-Set Gate           & 3.8  & 2.5 & 95.0 & 98.8 & 0.56/1.05 \\
    CommitGuard             & 2.5  & 1.3 & 93.8 & 98.8 & 0.51/1.03 \\
    \midrule
    \method{} w/o View Bind & 2.5  & 1.3 & 96.3 & 98.8 & 0.49/0.99 \\
    \method{} w/o Revalid.  & 87.5 & 1.3 & 98.8 & 100  & 0.49/0.92 \\
    \textbf{\method{} Gate} & \textbf{0} & \textbf{0} & \textbf{98.8} & \textbf{100} & 0.70/1.28 \\
    \bottomrule
    \end{tabular}%
    }
\end{table}
Table~\ref{tab:commit-results} evaluates commit-time dependency invalidation.
\method{} blocks unsafe commits triggered by all tested action-dependency and precondition mutations while preserving 98.8\% of valid actions under irrelevant changes.
With fixed proposals, removing view-to-action binding raises action-dependency UCR from 0 to 2.5\%; disabling live dependency revalidation raises it to 87.5\%, confirming that the safeguards are complementary.
\subsection{IPO Benchmark}
\begin{table}[tb]
    \centering
    \footnotesize
    \setlength{\tabcolsep}{4pt}
    \caption{IPO results (three trials/task): task-equivalent counts, per-trial Calls, and success-normalized Tok./VBWC.}
    \label{tab:ip-transfer}
    \resizebox{\columnwidth}{!}{%
    \begin{tabular}{@{}l rrc rr @{}}
    \toprule
    \textbf{Method}
    & \textbf{VBWC} & \textbf{Q-VBWC} & \textbf{Macro F1}
    & \textbf{Calls/trial} & \textbf{Tok./VBWC} \\
    \midrule
    ACON               & 69/80 & 64/80 & 0.816 & 6.0 & 42.1k \\
    ContextWeaver      & 72/80 & 68/80 & 0.854 & 8.9 & 48.7k \\
    ReadAgent-P        & 55/80 & 49/80 & 0.632 & 9.7 & 121.0k \\
    HiAgent            & 58/80 & 52/80 & 0.664 & 10.1 & 133.0k \\
    A-MEM              & 74/80 & 70/80 & 0.886 & 8.2 & 65.4k \\
    \midrule
    \textbf{\method{}}& \textbf{77/80} & \textbf{74/80} & \textbf{0.927} & \textbf{4.1} & \textbf{15.2k} \\
    \bottomrule
    \end{tabular}%
    }
\end{table}
On IPO, \method{} attains the highest task-macro VBWC (77/80; 96.3\%) and topology F1 (0.927), while using 15.2k tokens per successful workflow versus A-MEM's 65.4k (Table~\ref{tab:ip-transfer}).
These results indicate that bounded artifact projections and runtime resolution can limit context bloat and noise accumulation in long workflows.
\subsection{Public Benchmarks}
\begin{table}[tb]
    \centering
    \footnotesize
    \setlength{\tabcolsep}{4pt}
    \caption{Public transfer (three trials/task): task-equivalent Pass and trial-mean Token/Peak context.}
    \label{tab:public-transfer}
    \resizebox{\columnwidth}{!}{%
    \begin{tabular}{@{}l rrrc rrc @{}}
    \toprule
    \multirow{2}{*}{\textbf{Method}} & \multicolumn{4}{c}{\textbf{DeepSWE}} & \multicolumn{3}{c}{\textbf{Terminal-Bench}} \\
    \cmidrule(lr){2-5} \cmidrule(l){6-8}
    & \textbf{Pass} & \textbf{Part.} & \textbf{\shortstack{Tok.\\(M)}} & \textbf{\shortstack{Peak\\ctx.(k)}} & \textbf{Pass} & \textbf{\shortstack{Tok.\\(k)}} & \textbf{\shortstack{Peak\\ctx.(k)}} \\
    \midrule
    ACON               & 54/113 & 0.748 & 5.20 & 58.0 & 49/89 & 860 & 37.0 \\
    ContextWeaver      & 52/113 & 0.724 & 3.50 & 37.0 & 53/89 & 176 & 18.1 \\
    ReadAgent-P        & 48/113 & 0.701 & 4.20 & 50.5 & 46/89 & 920 & 22.5 \\
    HiAgent            & \textbf{60/113} & \textbf{0.825} & 2.60 & 32.8 & 51/89 & 710 & 28.5 \\
    A-MEM              & 51/113 & 0.713 & 2.80 & 28.5 & \textbf{64/89} & 262 & 14.7 \\
    \midrule
    \textbf{\methodCRI{}}& 56/113 & 0.766 & \textbf{1.10} & \textbf{14.7} & 55/89 & \textbf{128} & \textbf{9.3} \\
    \bottomrule
    \end{tabular}%
    }
\end{table}
Without custom Domain Contracts, \methodCRI{} evaluates only prospective context materialization and does not inherit Proposition~1's safety guarantee.
It attains task-equivalent Pass of 56/113 on DeepSWE and 55/89 on Terminal-Bench, with over 50\% fewer mean tokens and 36--55\% less peak context than the top-performing baselines (Table~\ref{tab:public-transfer}).
\section{Conclusion}
\label{sec:conclusion}
\method{} couples contract-admissible context construction with certificate-bound commit authorization to preserve reasoning-to-execution semantic continuity.
On governed tasks, it maintained evidence completeness; the full Commit Gate blocked all evaluated unsafe commits caused by dependency or precondition mutations; \methodCRI{} reduced public-task resource use.
Building on this verifiable foundation, our future work aims to generalize this semantic continuity paradigm to fully autonomous multi-agent ecosystems and unbounded dynamic environments.
%%%%%%%%%%%%%%%%%%%%%%%%%%%%%%%%%%%%%%%%%%%%%%%%%%%%%%%%%%%%%%%%%%%%%%%%
%%% The acknowledgments section is defined using the "acks" environment
%%% (rather than an unnumbered section). The use of this environment
%%% ensures the proper identification of the section in the article
%%% metadata as well as the consistent spelling of the heading.
%%%%%%%%%%%%%%%%%%%%%%%%%%%%%%%%%%%%%%%%%%%%%%%%%%%%%%%%%%%%%%%%%%%%%%%%
%%% The next two lines define, first, the bibliography style to be
%%% applied, and, second, the bibliography file to be used.
\balance
\bibliographystyle{ACM-Reference-Format}
\bibliography{arc_references}
\clearpage
\onecolumn
\appendix
\newcommand{\PromptSynopsis}[2]{%
\begin{center}
\begin{minipage}{0.92\textwidth}
\hrule
\par\strut{\large\bfseries #1}\par
{\large\ttfamily\raggedright #2\par}
\par\strut\hrule
\end{minipage}
\end{center}}
\section{Additional Method Details}
\label{app:method-details}
This appendix isolates the safety-critical implementation semantics underlying Section~\ref{sec:method}; deployment-specific logging and service configuration are omitted.
\subsection{Resolution, Certification, and Recovery}
\begin{table}[H]
\centering
\normalsize
\renewcommand{\arraystretch}{1.15}
\caption{Safety-critical division of responsibility.}
\begin{tabularx}{0.94\textwidth}{@{}p{0.19\textwidth} X X@{}}
\toprule
\textbf{Component} & \textbf{Operation} & \textbf{Integrity boundary} \\
\midrule
Requirement-state normalizer & Computes the normalized active plan $\widehat Q_s$. & Only committed declarations enter the active state; global requirements and pending declarations are runtime-protected. \\
Broker and compiler & Retrieve, rank, compress, and construct a candidate $V_s$. & Candidate-view generation only; they cannot certify a view. \\
Independent verifier & Recomputes Equation~\eqref{eq:admissible-view} from the contract, snapshot, manifest, budget, and canonical rendering. & A missing obligation or cross-snapshot substitution is rejected. \\
Recovery controller & Rebuilds before invocation or restarts after commit failure. & Pre-invocation rebuild retains $\sigma_s$; post-invocation failure discards $(a,Q^{\mathrm{decl}})$ and obtains a fresh snapshot. \\
\bottomrule
\end{tabularx}
\end{table}
The view selector uses the lexicographic priority
\[
\mathsf{mandatory\ feasibility}\succ\mathsf{optional\ coverage}\succ
\mathsf{rendering\ cost}\succ\mathsf{record\ count}\succ
\mathsf{canonical\ identifier}.
\]
Greedy and exact mixed-integer selection implement the same ordering; only Equation~\eqref{eq:admissible-view} determines acceptance.
\subsection{Certificate-Bound Sealing and Commit}
Algorithm~2 expands Lines 6--9 of Algorithm~1. A contract-version mismatch is verification failure, and proposal status is runtime-protected.
\begin{center}
\begin{minipage}{0.94\textwidth}
\hrule
\noindent{\large\textbf{Algorithm 2: Certificate-bound sealing and commit}}
\normalsize
\noindent\textbf{Input:} $\mathcal C_d$; $(\pi_s,V_s)$; CRI output $(\alpha,R_{\mathrm{add}},Q^{\mathrm{decl}})$ with $R_{\mathrm{add}}\subseteq\mathcal X$; reasoning snapshot $\sigma_s$.\\
\noindent\textbf{Output:} execution observation or $\bot$.
\renewcommand{\arraystretch}{1.10}
\begin{tabularx}{\linewidth}{@{}r@{\hspace{0.55em}}X@{}}
1: & Validate the CRI schema, guarded-resource identifiers, certificate integrity, and active contract version; otherwise return $\bot$ without creating a proposal. \\
2: & Freeze $D(a)\leftarrow\operatorname{Refs}(V_s)\cup\Delta_d(\alpha,\sigma_s)\cup\{(x,\nu_s(x)):x\in R_{\mathrm{add}}\}$. \\
3: & Seal one-shot $a$ with $\pi_s$, invocation identifier, $D(a)$, and pending $Q^{\mathrm{decl}}$; set $\operatorname{st}(a)\leftarrow\mathsf{pending}$. \\
4: & At the side-effect linearization point, obtain live $\sigma_e$ and enter one atomic conditional-execution block. \\
5: & If $\operatorname{st}(a)\ne\mathsf{pending}$, return $\bot$ without changing the terminal state. \\
6: & If verification, binding, freshness, or $P_d(\alpha,\sigma_e)$ fails, set $a$ to rejected, discard $Q^{\mathrm{decl}}$, and return $\bot$. \\
7a: & Attempt the managed application once. If it fails, roll back, set $a$ to rejected, discard $Q^{\mathrm{decl}}$, and return $\bot$. \\
7b: & On success, consume $a$, set it to committed, activate $Q^{\mathrm{decl}}$, and return the observation. \\
\end{tabularx}
\hrule
\end{minipage}
\end{center}
The required primitive is one atomic conditional transition spanning authorization, managed application, terminalization, and declaration activation. Non-participating external effects lie outside Proposition~1.
\section{Formal Analysis}
\label{app:formal-analysis}
\subsection{Operational Model}
For one governed decision, define the runtime configuration and proposal status by
\begin{equation}
\omega=\langle\sigma,\mathcal C_d,\pi,a,\mathcal Q\rangle,
\qquad
\ a\ne\varnothing\Longrightarrow\\
\operatorname{st}(a)\in
\{\mathsf{pending},\mathsf{committed},\mathsf{rejected}\},
\label{eq:runtime-state}
\end{equation}
where $\mathcal Q$ is the active requirement state and $\pi,a$ may be $\varnothing$. Write $z\triangleright\langle\cdot\rangle$ when a protected object $z$ binds the displayed fields. Then
\begin{align}
\pi_s &\triangleright
\langle \operatorname{ver}(\mathcal C_d),\iota_s,\sigma_s,
\widehat Q_s,B,(W_o)_{o\in O_d},V_s,
\operatorname{digest}(\operatorname{Render}(V_s))\rangle,               \label{eq:certificate-fields}\\
a &\triangleright
\langle\alpha,\pi_s,\iota_s,D(a),Q^{\mathrm{decl}},
\operatorname{st}(a)\rangle,                                            \label{eq:proposal-fields}
\end{align}
with invocation identifier $\iota_s$ and canonical rendering digest $\operatorname{digest}$. Let $\operatorname{Dep}^{\mathrm{gov}}_d(\alpha,\sigma)$ denote the set of contract-mandated resource--version pairs for $\alpha$ within the managed boundary.
\begin{table}[H]
\centering
\normalsize
\renewcommand{\arraystretch}{1.16}
\caption{Assumptions used by the safety proof.}
\begin{tabularx}{0.94\textwidth}{@{}p{0.07\textwidth} X X@{}}
\toprule
\textbf{ID} & \textbf{Formal condition} & \textbf{Scope} \\
\midrule
A1 & $x$ changes on $[\sigma,\sigma']\Rightarrow\nu_\sigma(x)\ne\nu_{\sigma'}(x)$. & Versions are fresh and non-reused; collection membership advances an aggregate version. \\
A2 & $(V,\sigma)\mapsto(\operatorname{Refs}(V),\operatorname{Render}(V),\operatorname{Cost}(\operatorname{Render}(V)))$ is deterministic. & Records are snapshot-immutable, $\operatorname{Refs}$ is transitively closed, and static input contains no omitted mutable managed-state premise. \\
A3 & $\operatorname{Verify}(\pi_s,\mathcal C_d)=1\Rightarrow\operatorname{Adm}_{\mathcal C_d}(V_s,\widehat Q_s,\sigma_s;B)=1$. & Verifier soundness; completeness is unnecessary. \\
A4 & $\operatorname{ver}(\pi_s)\ne\operatorname{ver}(\mathcal C_d)\Rightarrow\operatorname{Verify}(\pi_s,\mathcal C_d)=0$. & Equations~\eqref{eq:certificate-fields}--\eqref{eq:proposal-fields} are runtime-protected. \\
A5 & $\operatorname{Dep}^{\mathrm{gov}}_d(\alpha,\sigma)\subseteq\Delta_d(\alpha,\sigma)$; $P_d$ is type-correct. & Completeness is required only on the contract-governed managed boundary. \\
A6 & $\mathsf{authorize};\mathsf{apply};\mathsf{terminalize};\mathsf{activate}$ has one linearization point. & Failure preserves managed state and the active requirement state. \\
\bottomrule
\end{tabularx}
\end{table}
Let $\operatorname{Apply}(\alpha,\sigma_e)$ return the updated managed snapshot or $\bot$. A governed trace starts from
\[
\omega_0=\langle\sigma_s,\mathcal C_d,\varnothing,\varnothing,\mathcal Q\rangle .
\]
Between protocol transitions, environment evolution is represented by
\begin{equation*}
\textup{(E)}\qquad
\frac{\sigma\longrightarrow\sigma'\ \text{obeys A1}}
{\langle\sigma,\mathcal C_d,\pi,a,\mathcal Q\rangle
\xrightarrow{\textsc{Env}}
\langle\sigma',\mathcal C'_d,\pi,a,\mathcal Q\rangle},
\end{equation*}
where $\textsc{Env}$ cannot alter $\pi$, $a$, $\operatorname{st}(a)$, or $\mathcal Q$. A contract change is visible through $\mathcal C'_d$ and invalidates the old certificate by A4. The trusted runtime induces the following protocol transitions:
\begin{equation*}
\textup{(T1)}\qquad
\frac{\operatorname{Verify}(\pi_s,\mathcal C_d)=1}
{\langle\sigma_s,\mathcal C_d,\varnothing,\varnothing,\mathcal Q\rangle
\xrightarrow{\textsc{Certify}}
\langle\sigma_s,\mathcal C_d,\pi_s,\varnothing,\mathcal Q\rangle}
\end{equation*}
\begin{equation*}
\textup{(T2)}\qquad
\frac{\begin{gathered}
(\alpha,R_{\mathrm{add}},Q^{\mathrm{decl}})
=\operatorname{Agent}(C_s^{\mathrm{static}}\oplus\operatorname{Render}(V_s)),\\
\operatorname{Verify}(\pi_s,\mathcal C_d)=1,\qquad \mathsf{CRIValid}=1
\end{gathered}}
{\langle\sigma_s,\mathcal C_d,\pi_s,\varnothing,\mathcal Q\rangle
\xrightarrow{\textsc{Invoke--Seal}}
\langle\sigma_s,\mathcal C_d,\pi_s,a,\mathcal Q\rangle}
\end{equation*}
where
\begin{equation*}
\textup{(T2a)}\qquad
\operatorname{st}(a)=\mathsf{pending},\qquad
D(a)=\operatorname{Refs}(V_s)\cup\Delta_d(\alpha,\sigma_s)
\cup\{(x,\nu_s(x)):x\in R_{\mathrm{add}}\}.
\end{equation*}
\begin{equation*}
\textup{(T3)}\qquad
\frac{\begin{gathered}
\operatorname{st}(a)=\mathsf{pending},\quad
\Authorize_{\mathcal C_d}(a,\pi_s,\sigma_e)=1,\\
\operatorname{Apply}(\alpha,\sigma_e)=\sigma'_e
\end{gathered}}
{\langle\sigma_e,\mathcal C_d,\pi_s,a,\mathcal Q\rangle
\xrightarrow{\textsc{Commit}}
\langle\sigma'_e,\mathcal C_d,\pi_s,
a[\operatorname{st}:=\mathsf{committed}],
\operatorname{Activate}(\mathcal Q,Q^{\mathrm{decl}})\rangle}
\end{equation*}
\begin{equation*}
\textup{(T4)}\qquad
\frac{\begin{gathered}
\operatorname{st}(a)=\mathsf{pending},\\
\Authorize_{\mathcal C_d}(a,\pi_s,\sigma_e)=0
\ \lor
[\Authorize_{\mathcal C_d}(a,\pi_s,\sigma_e)=1
\land\operatorname{Apply}(\alpha,\sigma_e)=\bot]
\end{gathered}}
{\langle\sigma_e,\mathcal C_d,\pi_s,a,\mathcal Q\rangle
\xrightarrow{\textsc{Reject}}
\langle\sigma_e,\mathcal C_d,\pi_s,
a[\operatorname{st}:=\mathsf{rejected}],\mathcal Q\rangle}
\end{equation*}
No rule can apply $\alpha$ when
$\operatorname{st}(a)\in\{\mathsf{committed},\mathsf{rejected}\}$. A rejected proposal's pending $Q^{\mathrm{decl}}$ is permanently ineligible for activation.
\subsection{Safety Lemmas and Proposition}
\noindent\textbf{Lemma 1 (Certified-view origin).}
\begin{equation}
\omega\xrightarrow{\textsc{Invoke--Seal}}\omega'
\Longrightarrow
\operatorname{Adm}_{\mathcal C_d}(V_s,\widehat Q_s,\sigma_s;B)=1,
\quad C_s=C_s^{\mathrm{static}}\oplus\operatorname{Render}(V_s).
\label{eq:lemma-certified-origin}
\end{equation}
\emph{Proof.}
Rule~\textup{(T2)} requires verification and fixes the displayed model input. A3 gives admissibility; Equations~\eqref{eq:certificate-fields}--\eqref{eq:proposal-fields} and A2--A4 bind the admitted $V_s$ and its canonical rendering to $\iota_s$. \hfill$\square$
\noindent\textbf{Lemma 2 (Sealed-scope integrity).}
\begin{align}
\operatorname{Bound}_{\mathcal C_d}(a,\pi_s)=1
\Longrightarrow{}&
\operatorname{st}(a)=\mathsf{pending}
\land a[\pi]=\pi_s
\land a[\iota]=\pi_s[\iota] \notag\\
&\land D(a)=\operatorname{Refs}(V_s)\cup\Delta_d(\alpha,\sigma_s)
\cup\{(x,\nu_s(x)):x\in R_{\mathrm{add}}\}.
\label{eq:lemma-sealed-scope}
\end{align}
\emph{Proof.}
Rule~\textup{(T2)} constructs the displayed fields; A4 makes them immutable; $\operatorname{Bound}$ rejects every altered or non-pending proposal. \hfill$\square$
\noindent\textbf{Lemma 3 (Commit-point validity).}
\begin{equation}
\omega\xrightarrow{\textsc{Commit}}\omega'
\Longrightarrow
\left[\forall(x,\bar\nu)\in D(a):\nu_e(x)=\bar\nu\right]
\land P_d(\alpha,\sigma_e)=1.
\label{eq:lemma-commit-validity}
\end{equation}
\emph{Proof.}
By Rule~\textup{(T3)}, authorization equals one. Expanding Equation~\eqref{eq:commit-authorization} gives both conjuncts; A6 locates them at the application linearization point, and A1 excludes version reuse. \hfill$\square$
\noindent\textbf{Lemma 4 (One-shot fail-closed execution).}
For every sealed proposal $a$ and execution trace $\tau$, let
$N^\tau_{\textsc{Commit}}(a)$ count Rule~\textup{(T3)} transitions for $a$.
For
$\omega=\langle\sigma,\mathcal C_d,\pi,a,\mathcal Q\rangle$ and
$\omega'=\langle\sigma',\mathcal C'_d,\pi',a',\mathcal Q'\rangle$,
\begin{equation}
N^{\tau}_{\textsc{Commit}}(a)\le1,\qquad
\omega\xrightarrow{\textsc{Reject}}\omega'
\Longrightarrow \sigma'=\sigma\land\mathcal Q'=\mathcal Q.
\label{eq:lemma-one-shot}
\end{equation}
\emph{Proof.}
Rules~\textup{(T3)}--\textup{(T4)} permit only
$\mathsf{pending}\to\mathsf{committed}$ or
$\mathsf{pending}\to\mathsf{rejected}$; neither terminal state enables Rule~\textup{(T3)}. The second claim is A6 and Rule~\textup{(T4)}. \hfill$\square$
\noindent\textbf{Proof of Proposition 1.}
For an arbitrary successful commit, Rule~\textup{(T3)} and Equation~\eqref{eq:commit-authorization} imply
\begin{equation}
\begin{aligned}
&\operatorname{Verify}(\pi_s,\mathcal C_d)=1
\land\operatorname{Bound}_{\mathcal C_d}(a,\pi_s)=1\\
&\quad{}\land[\forall(x,\bar\nu)\in D(a):\nu_e(x)=\bar\nu]
\land P_d(\alpha,\sigma_e)=1.
\end{aligned}
\label{eq:prop-conjuncts}
\end{equation}
Rule~\textup{(T2)} is the only rule that creates a pending proposal, so every Rule~\textup{(T3)} transition has a preceding \textsc{Invoke--Seal} transition. Applying Lemmas~1--3 therefore yields
\begin{equation}
\begin{aligned}
&\operatorname{Adm}_{\mathcal C_d}(V_s,\widehat Q_s,\sigma_s;B)=1
\land a[\pi]=\pi_s
\land a[\iota]=\pi_s[\iota]\\
&\quad{}\land[\forall(x,\bar\nu)\in D(a):\nu_e(x)=\bar\nu]
\land P_d(\alpha,\sigma_e)=1,
\end{aligned}
\label{eq:prop-result}
\end{equation}
which is Proposition~1. Lemma~4 additionally gives one-shot, fail-closed execution, while A5 fixes the contract-relative dependency boundary. \hfill$\square$
Pre-invocation reconstruction and post-invocation recovery obey
\[
\frac{\operatorname{Verify}(\pi'_s,\mathcal C_d)=1}
{(\sigma_s,V_s,\pi_s)\xrightarrow{\mathsf{rebuild}}
(\sigma_s,V'_s,\pi'_s)},
\qquad
a_{\mathsf{rejected}}\not\xrightarrow{\textsc{Commit}},
\]
so reconstruction at fixed $\sigma_s$ preserves the invariant. Post-invocation rejection instead discards $(a,Q^{\mathrm{decl}})$, acquires a fresh reasoning snapshot, recomputes $\widehat Q$, and constructs a new view, certificate, and proposal. The result is a safety property, not a liveness or task-correctness guarantee.
\section{Internal Benchmarks and Sanitized Illustrations}
\label{app:benchmark-illustrations}
\subsection{Disclosure Boundary}
CDI and IPO are internally developed, non-public research benchmarks. The cases and prompt summaries below are purpose-built abstractions: no displayed constant, name, evidence fragment, or topology element is drawn from an evaluated task or used to compute the reported results. Prompt summaries are non-verbatim and non-executable; they preserve only task semantics and actor-visible information structure and do not constitute a benchmark release.
\subsection{Task Constructs and Protocols}
\begin{table}[H]
\centering
\normalsize
\renewcommand{\arraystretch}{1.16}
\caption{Internal benchmark constructs, abstract action protocols, and independent endpoint authorities.}
\label{tab:benchmark-constructs}
\begin{tabularx}{0.95\textwidth}{@{}p{0.13\textwidth} X X@{}}
\toprule
\textbf{Benchmark} & \textbf{Evaluated construct and actor protocol} & \textbf{Independent authority} \\
\midrule
CDI & Recover a source assignment for a versioned modular-affine DAG from a terminal target and progressively exposed evidence. At each decision, the actor either \textsc{Invokes} an admissible inverse operation to expand the reverse frontier or \textsc{Finishes} with a candidate and evidence lineage. $B_i^\star$ is the offline-oracle minimum budget sufficient for all mandatory evidence. & A deterministic interpreter executes the frozen graph forward and tests terminal equivalence. The private graph, source assignment, obligation labels, and mutation labels remain hidden. View completeness is evaluated separately by the blinded protocol in Section~\ref{app:metric-details}. \\
IPO & Recover canonical nodes, ports, and typed links in 80 physical-topology tasks. The persistent topology artifact evolves through
$\textsc{Preprocess}\to\textsc{Recover}\to\textsc{Readiness}\to\textsc{Publish}$;
each transition binds its immediate predecessor. & After terminal capture, a deterministic evaluator checks phase continuity, evidence grounding, canonical identifiers, typed links, and agreement with the private reference graph. Configuration-audit and risk-analysis profiles are outside the frozen evaluation. \\
\bottomrule
\end{tabularx}
\end{table}
\subsection{Purpose-Built Synthetic Illustrations}
\begin{table}[H]
\centering
\normalsize
\renewcommand{\arraystretch}{1.18}
\caption{Synthetic illustrations for interpretation; no displayed element occurs in an evaluated task.}
\label{tab:synthetic-examples}
\begin{tabularx}{0.95\textwidth}{@{}p{0.13\textwidth} X X@{}}
\toprule
\textbf{Benchmark} & \textbf{Reader-visible construction} & \textbf{Illustrative resolution} \\
\midrule
CDI & Over $\mathbb Z_{17}$, with the task-visible clue $x_0=2$,
$n_0=(x_0+3x_1+2)\bmod17$,
$n_1=(n_0+5x_0+4)\bmod17$, and $n_1=13$.
The complete equations are exposed only for this illustration. & A reverse trace yields $n_0=16$, $x_0=2$, and then $x_1=4$. Forward evaluation gives
$(2+3\cdot4+2)\bmod17=16$ and
$(16+5\cdot2+4)\bmod17=13$. \\
IPO & Two invented observations state that node $A$ at port $p_1$ is adjacent to node $B$ at port $p_9$, and conversely. & The supported artifact contains nodes $A,B$ and one typed link
$A{:}p_1\leftrightarrow B{:}p_9$, grounded in both observations. The private reference graph remains actor-invisible. \\
\bottomrule
\end{tabularx}
\end{table}
\PromptSynopsis{Sanitized, non-verbatim prompt synopses}{%
\textbf{CDI}\\
Objective: recover a source assignment consistent with the terminal observation under forward execution.\\
Evidence: ground each inverse step in currently visible evidence and preserve provenance.\\
Decision: invoke the next admissible inverse operation while incomplete; otherwise finish with the candidate.\\[0.7em]
\textbf{IPO}\\
Objective: recover the canonical physical topology from device and adjacency evidence.\\
Workflow: proceed from evidence preparation through recovery and readiness validation to publication.\\
Grounding: preserve predecessor continuity and ground every node and link in task-visible evidence.}
The synopsis specifies semantic intent only; it is neither the literal evaluation prompt nor a valid action payload.
\section{Experimental Protocol}
\label{app:experimental-protocol}
\subsection{Aggregation and Evaluation Runs}
All end-to-end comparisons use three trials per task. For trial-level outcome $m_{ir}$ on task $i$ and repeat $r$,
\begin{equation}
\overline m=\frac{1}{N}\sum_{i=1}^{N}\left(\frac{1}{3}\sum_{r=1}^{3}m_{ir}\right).
\label{eq:task-macro}
\end{equation}
\begin{table}[H]
\centering
\normalsize
\renewcommand{\arraystretch}{1.14}
\caption{Aggregation conventions.}
\begin{tabularx}{0.90\textwidth}{@{}p{0.22\textwidth} X@{}}
\toprule
\textbf{Convention} & \textbf{Definition} \\
\midrule
Analysis unit & Trial scores are averaged within task and then equally across the $N$ assigned tasks by Equation~\eqref{eq:task-macro}. \\
Failure handling & Refusal, fail-closed return, exhausted budget, provider/schema error, timeout, verifier error, or missing delivery receives zero; no replacement trial is added. \\
Count-style display & $N\overline m/N$ (e.g., $77/80$) is a task-equivalent aggregate, not a single-trial success count. \\
Special denominators & Obligation-, action-, and cost-normalized quantities use the denominators in Section~\ref{app:metric-details}. \\
\bottomrule
\end{tabularx}
\end{table}
\begin{table}[H]
\centering
\normalsize
\renewcommand{\arraystretch}{1.14}
\setlength{\tabcolsep}{5pt}
\caption{Evaluation actors, execution dates, and sample sizes.}
\label{tab:appendix-eval-overview}
\begin{tabularx}{0.88\textwidth}{@{}l X c@{}}
\toprule
\textbf{Benchmark} & \textbf{Actor and execution date} & \textbf{Tasks $\times$ trials} \\
\midrule
CDI & GLM-5.2-200K; Aug. 15--16, 2026 & $80\times3$ \\
IPO & GLM-5.2-200K; Aug. 22, 2026 & $80\times3$ \\
DeepSWE v1.1 & deepseek-v4-flash; Aug. 23--24, 2026 & $113\times3$ \\
Terminal-Bench 2.1 & deepseek-v4-flash; Aug. 23--24, 2026 & $89\times3$ \\
\bottomrule
\end{tabularx}
\end{table}
\subsection{Budgets and Execution Environment}
\begin{table}[H]
\centering
\normalsize
\renewcommand{\arraystretch}{1.14}
\caption{Actor-side evaluation ceilings. All calls use temperature zero, extended thinking disabled, and no provider-level retry.}
\begin{tabularx}{0.94\textwidth}{@{}p{0.15\textwidth} p{0.16\textwidth} p{0.18\textwidth} X@{}}
\toprule
\textbf{Setting} & \textbf{Actor turns} & \textbf{Active context} & \textbf{Output/call ceiling} \\
\midrule
CDI & Frozen task-specific & Frozen task-specific & Frozen evidence and call ceilings; structured output \\
IPO & $4$ & $56$k & $24{,}576$ cumulative actor-output tokens; structured output \\
Public benchmarks & $160$ & $64$k & $128$k cumulative assistant-output tokens; $16{,}384$ per response \\
\bottomrule
\end{tabularx}
\end{table}
\begin{table}[H]
\centering
\normalsize
\renewcommand{\arraystretch}{1.14}
\caption{Execution environment.}
\begin{tabularx}{0.94\textwidth}{@{}p{0.22\textwidth} X@{}}
\toprule
\textbf{Layer} & \textbf{Specification} \\
\midrule
Inference & Remote model serving; all provider-reported calls and tokens are charged to the originating method. \\
Local controller & Python 3.12, LangGraph 1.2.10, PuLP/CBC 3.3.2; dual Intel Xeon E5-2690 v3, 377 GiB RAM, no local inference GPU. \\
DeepSWE sandbox & Isolated Docker; 2 CPUs, 8 GiB RAM, 20 GiB storage, no agent-side network, and a separate verifier. \\
Terminal-Bench sandbox & Official per-task CPU, memory, storage, and timeout settings. \\
\bottomrule
\end{tabularx}
\end{table}
\subsection{Public-Benchmark Transfer}
\begin{table}[H]
\centering
\normalsize
\renewcommand{\arraystretch}{1.14}
\caption{Public-benchmark transfer boundary.}
\begin{tabularx}{0.94\textwidth}{@{}p{0.20\textwidth} p{0.20\textwidth} X@{}}
\toprule
\textbf{Benchmark} & \textbf{Frozen release} & \textbf{Evaluation boundary} \\
\midrule
DeepSWE & v1.1, 113 tasks & Official verifier; complete snapshot \\
Terminal-Bench & v2.1, 89 tasks & Official verifier; complete release \\
\midrule
\multicolumn{3}{@{}p{0.94\textwidth}@{}}{\methodCRI{} uses no custom Domain Contract or hidden witness annotation and therefore does not inherit Proposition~1. Reported values are three-trial task-macro estimates, not direct rankings against leaderboard protocols with different rollout counts.} \\
\bottomrule
\end{tabularx}
\end{table}
\subsection{Information Parity and Baseline Adaptation}
\begin{table}[H]
\centering
\normalsize
\renewcommand{\arraystretch}{1.14}
\caption{Information-parity constraints.}
\begin{tabularx}{0.94\textwidth}{@{}p{0.20\textwidth} X@{}}
\toprule
\textbf{Category} & \textbf{Constraint} \\
\midrule
Shared & Task statement, task-visible evidence, actor, tools and schemas, execution environment, actor-context ceiling, cumulative-token ceiling, and delivery schema. \\
Native representation & Identical domain semantics may be carried by a contract, skill, memory record, or system prompt according to the baseline mechanism. \\
Hidden from all actors & Answers, dependency and mutation labels, judge outputs, and verifier outputs. \\
Cost accounting & Method-native auxiliary ceilings follow each frozen call graph and are not transferable to actor calls; all auxiliary requests and tokens remain charged. \\
\bottomrule
\end{tabularx}
\end{table}
The baselines are mechanism-preserving common-harness adapters rather than author-maintained reference systems; only action schemas, artifacts, receipts, failure accounting, and terminal delivery are standardized.
\begin{table}[H]
\centering
\normalsize
\renewcommand{\arraystretch}{1.14}
\setlength{\tabcolsep}{5pt}
\caption{Mechanisms retained by the common-harness context adapters.}
\label{tab:appendix-adapters}
\begin{tabularx}{0.94\textwidth}{@{}l X@{}}
\toprule
\textbf{Method} & \textbf{Retained mechanism; commit policy} \\
\midrule
ACON & Bounded two-stage rolling-context compression; direct \\
ContextWeaver & Selective dependency-structured memory with local parent closure; direct \\
ReadAgent-P & Tool-result gisting and one original-page recall; direct \\
HiAgent & Current-subgoal detail with completed-subgoal summaries; direct \\
A-MEM & Task-local structured notes, links, and retrieval/evolution; direct \\
\method{} & Requirement closure with contract-admissible materialization; evidence-bound dependency gate \\
\bottomrule
\end{tabularx}
\end{table}
\section{Metric and Measurement Details}
\label{app:metric-details}
\subsection{CDI Evidence and Completion Metrics}
For task $i$, trial $r$, and governed decision $k$, let $\mathcal O^{\mathrm{eval}}_{irk}$ be the hidden evaluation-only obligation set, defined independently of the runtime contract obligations $O_d$. Let $J_{irk}(o)\in\{0,1\}$ be the blinded semantic verdict, $I_{irk}$ the actor-invocation indicator, and $L_{irk}$ the deterministic lineage-validity indicator. Then
\begin{equation}
\begin{aligned}
\operatorname{SC}_{irk}
&=\frac{1}{|\mathcal O^{\mathrm{eval}}_{irk}|}\sum_{o\in\mathcal O^{\mathrm{eval}}_{irk}}J_{irk}(o),\\
\operatorname{CV}_{irk}
&=\prod_{o\in\mathcal O^{\mathrm{eval}}_{irk}}J_{irk}(o),\\
\operatorname{CI}_{irk}
&=\operatorname{CV}_{irk}I_{irk}L_{irk}.
\end{aligned}
\label{eq:cdi-context-metrics}
\end{equation}
Let $K_{ir}$ contain every governed decision point required by the frozen protocol for one trial; if execution does not reach an actor invocation at such a point, $I_{irk}=0$. Let $S_{ir}$ indicate endpoint success and $A_{ir}$ indicate that no action labelled invalid by the frozen mutation oracle was committed. Then
\begin{equation}
\begin{aligned}
M_{ir}&=\frac{1}{|K_{ir}|}\sum_{k\in K_{ir}}M_{irk},
\quad M\in\{\operatorname{SC},\operatorname{CV},\operatorname{CI}\},\\
E_{ir}&=\prod_{k\in K_{ir}}\operatorname{CI}_{irk},\\
\operatorname{VBAC}
&=\frac{1}{N}\sum_{i=1}^{N}
\left(\frac{1}{3}\sum_{r=1}^{3}S_{ir}E_{ir}A_{ir}\right).
\end{aligned}
\label{eq:cdi-trial-aggregation}
\end{equation}
Equation~\eqref{eq:vbac} suppresses the pair $(i,r)$ with a flat episode index. Because every task has exactly three trials, that notation and Equation~\eqref{eq:cdi-trial-aggregation} yield the same estimate.
\begin{table}[H]
\centering
\normalsize
\renewcommand{\arraystretch}{1.14}
\caption{Blinded semantic-coverage protocol.}
\begin{tabularx}{0.94\textwidth}{@{}p{0.18\textwidth} X@{}}
\toprule
\textbf{Element} & \textbf{Definition} \\
\midrule
Obligations & Generated by the frozen private operator DAG and task-visible evidence semantics without method output. \\
Judge input & One obligation-specific semantic reference and the actor-visible view; deepseek-v4-flash, Aug. 15--16, 2026. \\
Withheld & Method identity, endpoint outcome, generated action, verifier feedback, and other obligation labels. \\
Verdict & $J_{irk}(o)=1$ iff all required semantic elements are explicit or entailed; deterministic code computes Equation~\eqref{eq:cdi-context-metrics}. \\
\bottomrule
\end{tabularx}
\end{table}
\begin{table}[H]
\centering
\normalsize
\renewcommand{\arraystretch}{1.14}
\caption{CDI metric authority and denominator.}
\begin{tabularx}{0.95\textwidth}{@{}p{0.20\textwidth} p{0.29\textwidth} X@{}}
\toprule
\textbf{Metric} & \textbf{Positive condition / denominator} & \textbf{Authority} \\
\midrule
Semantic Coverage & Mean obligation-level semantic completeness over $\mathcal O^{\mathrm{eval}}_{irk}$ & Blinded semantic judge + deterministic aggregation \\
Complete View & $\operatorname{CV}_{irk}=1$ & Same as above \\
Complete Invocation & $\operatorname{CI}_{irk}=1$; does not require a correct action & View verdict + invocation log + lineage checker \\
Endpoint Success & Correct terminal candidate & Independent symbolic domain verifier \\
VBAC & Equation~\eqref{eq:cdi-trial-aggregation}, with $E_{ir}=\prod_{k\in K_{ir}}\operatorname{CI}_{irk}$ & Context protocol + endpoint verifier + mutation oracle \\
CVR & Committed governed actions with missing dependency; undefined if no governed action commits & Frozen dependency oracle \\
\bottomrule
\end{tabularx}
\end{table}
\subsection{IPO and Public-Benchmark Metrics}
\begin{table}[H]
\centering
\normalsize
\renewcommand{\arraystretch}{1.14}
\caption{IPO workflow and topology metrics.}
\begin{tabularx}{0.94\textwidth}{@{}p{0.18\textwidth} X@{}}
\toprule
\textbf{Metric} & \textbf{Definition} \\
\midrule
VBWC & Completed workflow, intact predecessor commitments, terminal publication, and valid evidence citations. \\
Q-VBWC & VBWC plus edge precision, edge recall, and evidence coverage each at least $0.95$. \\
Macro F1 & Equal-weight task macro-average of per-trial set F1 over canonical node identifiers and typed link keys. \\
\bottomrule
\end{tabularx}
\end{table}
For IPO task $i$ and trial $r$, let $T_i$ and $\widehat T_{ir}$ be the reference and predicted sets of typed topology elements, comprising canonical node identifiers and link keys. The trial score and task-macro aggregate are
\begin{equation}
F_{ir}=\frac{2\lvert T_i\cap\widehat T_{ir}\rvert}{\lvert T_i\rvert+\lvert\widehat T_{ir}\rvert},
\qquad
\operatorname{Macro\text{-}F1}=\frac{1}{N}\sum_{i=1}^{N}\left(\frac{1}{3}\sum_{r=1}^{3}F_{ir}\right).
\label{eq:ipo-macro-f1}
\end{equation}
A missing or invalid delivery, or a zero denominator, receives $F_{ir}=0$; ``Macro'' therefore denotes equal task weighting rather than class-wise averaging.
\begin{table}[H]
\centering
\normalsize
\renewcommand{\arraystretch}{1.14}
\caption{Public-benchmark and resource metrics.}
\begin{tabularx}{0.95\textwidth}{@{}p{0.18\textwidth} X X@{}}
\toprule
\textbf{Metric} & \textbf{Aggregation} & \textbf{Boundary} \\
\midrule
Pass / Partial & Trial mean within task, then task macro-average by Equation~\eqref{eq:task-macro} & Official verifier; Partial only for DeepSWE \\
Calls & Actor and method-native auxiliary requests per assigned trial & Includes failed trials \\
Token & Provider-reported input + output per assigned trial & Includes failed trials and auxiliary calls \\
Peak Context & Within-trial maximum active actor input, averaged across trials & Actor input only \\
Tok./VBWC & Total actor + auxiliary tokens divided by successful VBWC workflows & Undefined when no VBWC succeeds \\
\bottomrule
\end{tabularx}
\end{table}
\subsection{Mutation and Latency Protocol}
\begin{table}[H]
\centering
\normalsize
\renewcommand{\arraystretch}{1.14}
\caption{Commit-evaluation mutation strata. Every policy receives identical frozen proposals and traces; validity is assigned by the frozen benchmark oracle independently of the evaluated policy's runtime predicates.}
\begin{tabularx}{0.94\textwidth}{@{}p{0.22\textwidth} X p{0.15\textwidth}@{}}
\toprule
\textbf{Stratum} & \textbf{Perturbation} & \textbf{Oracle class} \\
\midrule
Action dependency & Mutate a resource labelled action-relevant by the frozen dependency oracle & Invalid \\
Precondition & Inject a state change labelled by the frozen benchmark oracle as violating an action precondition & Invalid \\
Irrelevant & Advance versions only for resources labelled action-irrelevant by the same oracle & Valid \\
No change & Preserve the evaluated state & Valid \\
\bottomrule
\end{tabularx}
\end{table}
Here, ``valid'' and ``invalid'' in Equation~\eqref{eq:commit-metrics} are benchmark-oracle labels. The Precondition stratum does not assert $P_d(\alpha,\sigma_e)=0$ for every evaluated policy: $P_d$ is the active contract predicate, whereas the mutation label is assigned independently from the frozen reference semantics. Consequently, a policy may accept an oracle-invalid case while its runtime predicate evaluates to one; Proposition~1 remains contract-relative. Latency is measured as follows:
\begin{table}[H]
\centering
\normalsize
\renewcommand{\arraystretch}{1.14}
\caption{Commit-gate latency boundary.}
\begin{tabularx}{0.94\textwidth}{@{}p{0.18\textwidth} X@{}}
\toprule
\textbf{Element} & \textbf{Measurement rule} \\
\midrule
Included & One in-process gate predicate, timed by a high-resolution monotonic clock. \\
Excluded & Model inference, network transport, snapshot acquisition, view compilation, database transactions, environment effects, and official verification. \\
Aggregation & p50 and p95 over all gate calls for a policy in one execution environment; interpreted as relative gate overhead, not end-to-end latency. \\
\bottomrule
\end{tabularx}
\end{table}
\end{document}
